\documentclass[11pt]{report}
\usepackage[margin=1in]{geometry}
\usepackage{graphicx,amsmath,amssymb,booktabs,longtable,hyperref,xcolor,tabularx,float}
\usepackage{caption}
\graphicspath{{./}}
\hypersetup{colorlinks=true,linkcolor=blue,urlcolor=blue,citecolor=blue}
\setlength{\parskip}{5pt}
\setcounter{tocdepth}{2}
\newcommand{\Fw}{F_{\!w}}
\newcommand{\Lpriv}{L_{\rm priv}}
\newcommand{\Lpub}{L_{\rm pub}}
\newcommand{\strict}{strict-$h{=}0$}

\title{\textbf{Computing the K=3 CRRA Rational-Expectations Equilibrium:}\\
A Complete Analysis of the Fixed-Point Factory Investigation\\[1ex]
\large Economic problem, numerical operators, seventeen findings,\\
and the recommended reproducible \strict{} gold-standard pipeline}
\author{Fixed-Point Factory --- consolidated monograph}
\date{June 2026}

\begin{document}
\maketitle

\begin{abstract}
\noindent
This monograph consolidates the complete investigation of the K=3 CRRA
rational-expectations equilibrium (REE) with binary state and Gaussian
private signals: the economic question (does a deterministic, noiseless
CRRA economy support a \emph{partially revealing} equilibrium?), the
numerical machinery built to answer it (seven operator families, five
precision regimes, three coordinate systems), the seventeen numbered
findings produced along the way --- both positive and negative --- and an
extensive, step-by-step recommendation for the reproducible
\strict{} computation that is closest to the truth.

The headline economics: \textbf{the deterministic K=3 CRRA economy has a
genuine, smooth, partially revealing equilibrium --- no noise traders
needed.} The price equals the risk-weighted average belief plus a
Jensen/curvature wedge $(\tfrac12 - p)\,\mathrm{Var}_W(m)/\gamma$ that
is identically zero under full revelation and under CARA preferences.
The headline numerics: the \emph{lookup side} of the fixed-point problem
(the conditional posterior $\mu(p, u_k)$) is solved to machine epsilon by
a per-segment Chebyshev construction with cusp-located knots, across the
entire $(\gamma,\tau)$ grid tested; the \emph{equilibrium-surface side}
is nailed to $10^{-14}$ at low signal precision and to $\sim 10^{-12}$
with the consistent kernel co-area operator generally, while the bare
co-area operator at high resolution is provably $C^0$-but-not-$C^1$ at
Morse-critical prices --- a numerical fact about the operator, not an
economic statement about the equilibrium.
\end{abstract}

\tableofcontents
\listoffigures

% =====================================================================
\chapter*{Executive summary}
\addcontentsline{toc}{chapter}{Executive summary}
% =====================================================================

\section*{The question}
Since Grossman--Stiglitz (1980) and Hellwig (1980), the standard route to
a partially revealing asset price has been exogenous noise: noise
traders, random supply, or signal noise that prevents the price from
fully aggregating private information. The question pursued here is
whether a fully \emph{deterministic} economy --- $K$ CRRA traders, a binary
asset value $v \in \{0,1\}$, Gaussian private signals, no noise of any
kind --- can have an equilibrium price that does \emph{not} fully reveal
the pooled information.

\section*{The answer}
Yes, for $K = 3$. The deterministic K=3 CRRA economy has a genuine,
smooth, partially revealing (PR) equilibrium, computable to residual
$10^{-12}$--$10^{-14}$, with a positive, grid-independent information
deficit. The mechanism is the \emph{Jensen gap}: expanding CRRA demand in
log-odds ($m = \operatorname{logit}\mu$, $\pi = \operatorname{logit}p$,
$\delta = (m - \pi)/\gamma$) gives $x_i \approx W_i[\delta_i +
(\tfrac12 - p)\delta_i^2]$, and market clearing yields, to second order,
\[
  \pi \;=\; \bar m_W \;+\; \Bigl(\tfrac12 - p\Bigr)\,
            \frac{\mathrm{Var}_W(m)}{\gamma},
\]
price = average belief + a curvature wedge proportional to belief
\emph{disagreement}. The wedge vanishes under full revelation (no
disagreement) and under CARA ($\gamma \to \infty$ limit; linear demand,
no curvature). It compresses the price toward $\tfrac12$ ---
under-reaction, measured slope $\beta < \tau$. K=4 destroys the effect
through aggregation (disagreement $\to$ 0 $\Rightarrow$ full revelation
is globally attracting), unless extreme risk-aversion heterogeneity
sidelines agents and resurrects PR through an effective-agent-count
channel.

\section*{The numerical journey, compressed}
\begin{enumerate}\itemsep2pt
  \item \textbf{The original sin: the naive contour scan.} The
        equilibrium operator requires the level-set (co-area) integral
        $A_v(p) = \int_{\{P = p\}} f_v / |\nabla P|\, d\sigma$. The
        naive grid scan is inconsistent twice over: binary edge
        inclusion creates $O(0.1)$ discontinuities, and the missing
        $1/|\nabla P|$ weight biases the measure by $\sim$19\%. Years
        of apparent ``non-smooth residual floors'' were this artifact.
  \item \textbf{The kernel as consistent quadrature.} A Gaussian band
        $\int K_h(P - p) f\, du$ converges to the correct co-area
        integral as $h \to 0$ (verified to $1.5 \times 10^{-5}$); taken
        in the joint limit $\Delta x \ll h \to 0$, it nails the PR
        equilibrium to $10^{-14}$. The kernel is a numerical
        convergence knob, \emph{not} hidden economic noise: the deficit
        is invariant to $h$ once resolved.
  \item \textbf{But fixed-$h$ kernel-band extrapolation fails.} The
        kernel-band Richardson (R2--R7) lookup on a fixed grid does
        \emph{not} converge to \strict{}: level-set kinks inject
        non-integer powers of $h$ that even-power Richardson cannot
        annihilate. Lookup bias up to 0.6 in $\mu$; the K=3
        ``deficit ridge'' reported on that operator was an artifact
        (36-cell recheck: deficit is flat in $\tau$ at $\sim 10^{-2}$
        under \strict{}, with a sign-flipping bias crossing at
        $\tau \approx 0.4$).
  \item \textbf{The lookup problem is solved.} Detect the cusps of
        $\mu(p, u_k)$ from $\mu$ itself (the 3D critical points of $P$
        are \emph{not} the right knots --- slice criticality is), clean
        the no-roots fallback with Bayesian PCHIP extrapolation, then
        fit degree-6 Chebyshev per segment between cusp knots:
        machine-$\varepsilon$ median across all 36 $(\gamma, \tau)$
        cells, 11/11 monotone slices, correct boundary conditions, in
        7 seconds total.
  \item \textbf{The surface problem at high resolution is intrinsically
        hard.} The bare co-area evidence is $C^0$-but-not-$C^1$ at
        Morse-critical prices (proven on a controlled saddle: $A_v'(p)$
        blows up $128\times$ while the kernel-band derivative stays
        bounded). Newton needs $C^1$. Smoothed indicators introduce
        bias with no sweet spot; Newton--Krylov, MINPACK HYBRD, and
        frozen-cusp variants all plateau at the same
        $\sim 10^{-2}$ wall on $G = 11$ at high $\tau$. The fix is
        resolution + consistent kernel (the joint central-path limit),
        not solver engineering.
  \item \textbf{Multiplicity is real.} $P \equiv \tfrac12$ is an exact
        fixed point of the K=3 CRRA clearing at every $(\gamma, \tau)$;
        cold-started Newton finds it instantly. The informative PR
        branch requires warm-starting along a continuation ladder. The
        K=2 double-double work locates a fold at $\tau \approx 0.79$
        ($\sigma_{\min}(I - J)$: $0.45 \to 0.31 \to 0.017$ for
        $\tau = 1.0 \to 0.9 \to 0.8$) bounding its branch.
\end{enumerate}

\section*{The recommendation, compressed}
Part~VII gives the full protocol. In one paragraph: \emph{represent} the
equilibrium on a private-log-odds grid ($\xi = \tanh(\tau u/2)$ uniform,
no hard cube truncation, $\pm 5\sigma$ coverage); \emph{solve} with the
consistent kernel co-area operator in the joint limit
$\Delta x \ll h \to 0$ (the only operator that is simultaneously
consistent and $C^\infty$), warm-starting along a $\tau$-continuation
ladder from the analytic low-$\tau$ solution, with permutation and
sign-flip symmetry projection every iteration and a density-weighted
residual as the acceptance metric; \emph{extract} the lookup
$\mu(p, u_k)$ with the per-segment Chebyshev pipeline (v5) at degree 6;
\emph{verify} with the five-test battery (symmetry, monotonicity,
boundary conditions, operator cross-check against \strict{} co-area at
matched tolerance, and deficit-measure invariance); and \emph{report}
the deficit with its measure convention, the kernel-ladder evidence of
$h$-invariance, and the residual floor with its provenance. Everything
needed is committed on branch
\texttt{claude/study-fixed-point-economics-y12PB}.

% =====================================================================
\part{The economic problem}
% =====================================================================

\chapter{Rational expectations and the revelation question}

\section{The classical tension}
A competitive asset price aggregates dispersed private information. If
it aggregates \emph{perfectly} --- full revelation (FR) --- then no
trader has an incentive to acquire costly information, and the
equilibrium concept undermines itself (the Grossman--Stiglitz paradox).
The classical resolution adds exogenous randomness: noise traders,
stochastic asset supply, or unlearnable taste shocks. With noise in the
price, observation of $p$ leaves residual uncertainty about $v$, prices
are partially revealing (PR), and information acquisition can be
rewarded in equilibrium.

The price of that resolution is conceptual: the noise is an
\emph{assumption}, not a derivation. The investigation documented here
asks the sharper question: in a completely deterministic economy ---
no noise traders, no random supply, nothing unlearnable --- can the
equilibrium price still fail to reveal pooled information?

\section{Why determinism seems to force full revelation}
With finitely many signal profiles mapping to prices, a one-to-one price
function would let every agent invert the price into the pooled signal
vector (up to permutation, which is payoff-irrelevant under symmetric
weights). The only way a deterministic price can hide information is to
be \emph{many-to-one}: distinct signal profiles with distinct pooled
implications must map to the same price. The question is whether such a
many-to-one price can simultaneously clear the market at every signal
profile when every trader conditions on it rationally --- a fixed-point
requirement that is far from obviously satisfiable.

\section{The answer in brief}
For K=3 CRRA traders it is satisfiable. The equilibrium price function
$P(u_1, u_2, u_3)$ is smooth and genuinely many-to-one: level sets of
$P$ pool signal profiles with different posterior implications, and the
market still clears identically. The wedge that makes this possible is
the curvature (wealth-effect) term in CRRA demand --- absent under CARA,
which is why the CARA benchmark fully reveals (the Hellwig closed form
$p^* = \sigma(\tau \sum_i u_i)$ is the unique CARA equilibrium, verified
numerically to deficit $\sim 10^{-4}$, shrinking with grid).

\chapter{The model}

\section{Primitives}
\begin{itemize}\itemsep2pt
  \item Binary state $v \in \{0, 1\}$, symmetric prior
        $\mathbb{P}[v{=}1] = \tfrac12$.
  \item $K$ traders, $i = 1..K$, with CRRA utility of terminal wealth,
        risk-aversion parameter $\gamma$ and market weights $W_i$
        (homogeneous $W_i = 1$ unless stated).
  \item Private signals $u_i \mid v \sim \mathcal{N}(v - \tfrac12,\,
        \tau^{-1})$: precision $\tau$ controls informativeness. At
        $\tau \to 0$ signals are uninformative; at $\tau \gtrsim 1$
        a single signal is strongly informative
        ($\sigma_{\rm signal} = 1/\sqrt{\tau}$ vs the means gap of 1).
  \item One risky asset paying $v$, zero net supply; a numeraire bond.
\end{itemize}

\section{Information and demands}
Trader $i$ observes $(u_i, p)$ and forms the posterior
$\mu_i = \mathbb{P}[v{=}1 \mid u_i, P(\mathbf{u}) = p]$. With binary
payoff, the CRRA demand has the log-odds form: writing
$m_i = \operatorname{logit}\mu_i$, $\pi = \operatorname{logit}p$ and
$R_i = \exp((m_i - \pi)/\gamma)$, individual asset demand solves
\[
  x_i(p) \;=\; W_i\,\frac{R_i - 1}{(1 - p) + R_i\,p},
\]
which is the exact expression used in the production clearing solver
(bisection on $p$; \texttt{dd\_k3\_ops.py:dd\_crra\_clear}).
\textbf{A cautionary tale from this session}: the linearized
($\gamma \to \infty$) clearing $p = \frac1K \sum_i \mu_i$ looks similar
and is catastrophically wrong for the question at hand --- it removes
$\gamma$ from the fixed-point problem entirely, which an early run
exposed as the absurd result ``every $\gamma$ gives the identical
equilibrium''. The exact CRRA clearing must be used.

\section{Equilibrium as a fixed point}
A rational-expectations equilibrium is a price function
$P : \mathbb{R}^K \to (0,1)$ such that for every signal profile
$\mathbf{u}$,
\[
  \sum_{i=1}^{K} W_i\,
  \frac{R_i(\mathbf{u}) - 1}{(1 - P(\mathbf{u})) + R_i(\mathbf{u})\,P(\mathbf{u})}
  \;=\; 0,
  \qquad
  R_i = \exp\!\Bigl(\tfrac{\operatorname{logit}\mu_i - \operatorname{logit}P(\mathbf{u})}{\gamma}\Bigr),
\]
where each $\mu_i$ is the \emph{Bayes-correct} posterior given $u_i$ and
the event $\{P(\cdot) = P(\mathbf{u})\}$ computed under the same
function $P$. Operationally this is the loop
\[
  P \;\xrightarrow{\;\text{co-area integral}\;}\; \mu(p, u_k)
  \;\xrightarrow{\;\text{CRRA clearing}\;}\; \Phi(P),
  \qquad \text{REE} \iff \Phi(P) = P .
\]

\section{The lookup function}
The computational heart is the conditional posterior --- the
\emph{lookup function}
\[
  \mu(p, u_k) \;=\;
  \frac{ f_1(u_k)\, A_1(p, u_k) }
       { f_0(u_k)\, A_0(p, u_k) + f_1(u_k)\, A_1(p, u_k) },
  \qquad
  A_w(p, u_k) \;=\; \int_{\{P(\cdot, u_k) = p\}}
      \frac{f_w(u_a)\, f_w(u_b)}{|\nabla P|}\, d\sigma ,
\]
a Gaussian-weighted co-area integral over the level set of the
\emph{slice} of $P$ at the trader's own signal $u_k$, where
$f_w(u) = \mathcal{N}(u;\, w - \tfrac12, \tau^{-1})$. All seventeen
findings of Parts III--VI are about how to compute $A_w$ and the fixed
point accurately.

\section{Symmetries}
Four structural symmetries, used throughout both to halve unknowns and
as free correctness tests:
\begin{enumerate}\itemsep2pt
  \item \textbf{Permutation}: $P(u_{\pi(1)}, u_{\pi(2)}, u_{\pi(3)}) =
        P(\mathbf{u})$; reduces $G^3$ cube unknowns to
        $\binom{G+2}{3}$ (e.g.\ $1331 \to 286$ at $G{=}11$).
  \item \textbf{Sign flip}: $P(\mathbf{u}) + P(-\mathbf{u}) = 1$ (state
        relabeling); pins $P(\mathbf{0}) = \tfrac12$.
  \item \textbf{Axis monotonicity}: $\partial P/\partial u_i \ge 0$; a
        violation flags an operator artifact (the kernel-band R4 FPs
        violate it in 27\% of cells; \strict{} in 4\%; the structural
        ansatz in 0\%).
  \item \textbf{Lookup mirror}: $\mu(1 - p, -u_k) = 1 - \mu(p, u_k)$;
        any asymmetry of the lookup about $p = \tfrac12$ measures
        operator bias, not economics.
\end{enumerate}

\chapter{The Jensen gap: mechanism and closed form}

\section{Second-order expansion}
Expanding the CRRA demand in $\delta_i = (m_i - \pi)/\gamma$:
\[
  x_i \;\approx\; W_i\Bigl[\, \delta_i + \bigl(\tfrac12 - p\bigr)\,\delta_i^2 \,\Bigr].
\]
The quadratic term is the wealth-curvature effect, identically absent
for CARA. Imposing $\sum_i x_i = 0$ and solving to second order:
\[
  \boxed{\;\pi \;=\; \bar m_W \;+\;
  \Bigl(\tfrac12 - p\Bigr)\,\frac{\mathrm{Var}_W(m)}{\gamma}\;}
\]
the equilibrium log-odds price equals the weighted-average belief plus a
wedge proportional to belief \emph{dispersion} and inversely proportional
to risk aversion.

\section{Verified consequences}
All checked numerically on nailed equilibria (\texttt{k3\_jensen\_gap}):
\begin{itemize}\itemsep2pt
  \item \textbf{Gap $\equiv$ partial revelation.} Under FR all agents
        share the pooled posterior, $\mathrm{Var}_W(m) = 0$, the wedge
        vanishes --- so FR is always a clean fixed point, and any gap
        \emph{is} partial revelation.
  \item \textbf{Compression toward $\tfrac12$.} The sign factor
        $(\tfrac12 - p)$ pulls the price toward even odds:
        under-reaction; measured regression slope $\beta < \tau$.
  \item \textbf{$\propto 1/\gamma$}: more risk aversion $\Rightarrow$
        smaller wedge $\Rightarrow$ more revealing. CARA
        ($\gamma \to \infty$) is the no-gap benchmark; verified
        deficit $\sim \gamma^{-0.94} \to 0$ (\texttt{k3\_cara}).
  \item \textbf{$\propto \mathrm{Var}(m)$}: sharper signals (larger
        $\tau$) disperse beliefs and widen the wedge --- the deficit
        rises with $\tau$ and falls with $\gamma$ across the definitive
        $(\gamma, \tau)$ map (\texttt{k3\_coarea\_2dsweep}, $G{=}17$,
        31/32 cells nailed to $10^{-12}$).
  \item \textbf{K=4 collapse}: more agents aggregate beliefs, shrink
        disagreement, kill the wedge --- FR is unique and globally
        attracting at K=4 under homogeneous parameters (every seed
        tested flows to FR; \texttt{k4\_sym}, \texttt{k4\_basin},
        \texttt{k4\_extreme}).
  \item \textbf{Effective-agent-count channel}: with extreme
        heterogeneity $\gamma = [g, g, 1, 1]$, $g \to \infty$, the two
        risk-averse traders stop trading, the price reveals only the
        active pair, and genuine PR reappears (deficit $\to \sim 0.33$;
        \texttt{k4\_riskaversion}).
\end{itemize}

\section{The CARA/Hellwig cross-check}
The K=3 noiseless CARA model admits the sufficient statistic
$\log[f_1/f_0] = \tau \sum_i u_i$, so the unique CARA equilibrium is the
analytic fully-revealing price $p^* = \sigma(\tau \sum_i u_i)$. Direct
numerical evaluation at $P_{\rm FR}$ gives deficit $\sim 10^{-4}$,
shrinking with $G$ --- the Hellwig closed form is reproduced, and an
earlier ``growing CARA deficit'' was traced to a self-inconsistent
boundary halo, not economics
(\texttt{k3\_cara/HELLWIG\_CLOSED\_FORM.md}).

\chapter{What ``closest to the truth'' must mean}

Because the investigation produced both genuine economics and a long
catalogue of numerical artifacts that \emph{looked} like economics, the
monograph adopts an explicit standard for ``truth-closeness'' that
Part~VII operationalizes:

\begin{enumerate}\itemsep2pt
  \item \textbf{Consistency}: the discrete operator must converge to
        the continuum co-area operator as resolution increases ---
        including the $1/|\nabla P|$ weight (the naive scan's 19\% bias
        is disqualifying) and without discrete include/exclude flips.
  \item \textbf{Parameter-invariance of economics}: any auxiliary
        numerical parameter (kernel bandwidth, grid, quadrature order)
        must leave economic quantities (deficit, slope) invariant once
        resolved. The kernel passes this test; the fixed-grid
        Richardson extrapolation does not.
  \item \textbf{Symmetry exactness}: permutation, sign-flip, and lookup
        mirror must hold to within the arithmetic precision, not the
        truncation error.
  \item \textbf{Branch identification}: with multiple fixed points
        (trivial $P \equiv \tfrac12$, PR branch, FR), the computation
        must demonstrate \emph{which} branch it is on and how it got
        there (continuation provenance).
  \item \textbf{Honest residual reporting}: best-iterate semantics,
        stall amplitudes, and measure conventions for the deficit must
        be stated --- the difference between ``$|F| = 10^{-14}$'' and
        ``best iterate of an oscillating sequence with amplitude
        $0.05$'' decided several findings.
\end{enumerate}

% =====================================================================
\part{Numerical foundations}
% =====================================================================

\chapter{Why this problem is numerically hard}

\section{Three compounding difficulties}
\begin{enumerate}\itemsep2pt
  \item \textbf{The lookup is a level-set integral.} $A_w(p, u_k)$
        integrates over $\{P = p\}$ with weight $1/|\nabla P|$. Level
        sets change topology with $p$, hit domain boundaries, and the
        weight is singular at critical points of $P$. Each of these
        produces a kink or cusp in $\mu(p, u_k)$ that downstream
        numerics must represent without ringing.
  \item \textbf{The problem is self-consistent.} Lookup errors feed
        into the price via clearing, which feeds back into the lookup.
        An operator bias of 0.05 in $\mu$ --- invisible to the
        operator's own residual --- contaminates every economic
        quantity built on the fixed point.
  \item \textbf{Self-residual is not correctness.} The kernel-band R4
        operator converged to $|F| \sim 10^{-14}$ on its own fixed
        point while sitting $\sim 0.05$--$0.6$ away from the true
        lookup. The fixed point migrates to compensate for operator
        bias. The only defense is an independently derived reference
        operator --- the role \strict{} plays throughout Parts III--VI.
\end{enumerate}

\section{The co-area correction (the foundational artifact)}
The single most consequential numerical fact of the whole investigation:
the continuum inference $A_v(p) = \int_{\{P=p\}} f_v / |\nabla P|\,
d\sigma$ is a \emph{smooth} function of $p$, but the naive grid contour
scan (detect sign changes of $P - p$ on grid edges, sum $f_v$ over
detected edges) is inconsistent in two separately verified ways:
\begin{enumerate}\itemsep2pt
  \item \textbf{Binary include/exclude}: as $p$ crosses a grid node
        value, edges flip in or out discretely, producing $O(0.1)$
        jumps in the evidence that do \emph{not} shrink as the price
        step halves (\texttt{k3\_strict\_h0\_exact}).
  \item \textbf{Missing $1/|\nabla P|$}: the naive scan integrates
        arc length, not the co-area measure. On a known analytic
        surface: naive $= 2.07$ vs true $= 1.748$ --- a 19\% bias
        (\texttt{k3\_verify\_coarea}).
\end{enumerate}
Both are removed by either (a) exact marching-squares contours with the
explicit $1/|\nabla P|$ weight (the \strict{} operator), or (b) the
Gaussian kernel band $\int K_h(P - p) f\, du$, which converges to the
correct co-area integral as $h \to 0$ (verified to $1.5\times10^{-5}$).
For most of the project's pre-history, the apparent ``non-smooth
residual floor'' of the strict operator --- interpreted at the time as
``PR needs noise'' --- was exactly this artifact.

\section{$C^0$-but-not-$C^1$: the smoothness barrier, proven}
The deeper fact: even the \emph{exact} co-area evidence is only $C^0$ in
$p$ at Morse-critical prices (values $p_c = P(u^*)$ where
$\nabla P(u^*) = 0$ on the level set). On a controlled saddle surface,
$A_v(p)$ is continuous but its derivative blows up $128\times$ at the
critical price, while the kernel-band derivative stays bounded
(\texttt{k3\_hfree\_morse/c0\_not\_c1}). Newton-type solvers need
$C^1$; at finer grids more cells sit near critical prices; hence the
bare co-area operator is fundamentally not Newton-nailable at high
resolution. This explains, in one stroke, the entire family of
high-$\tau$ plateaus met in Parts III and VI: \strict{} Anderson stall
(Finding 9), zero-h solver plateau (Finding 17), Newton--Krylov and
HYBRD plateaus (Part VI).

\begin{figure}[H]\centering
  \includegraphics[width=0.8\linewidth]{cusps_figs/fig2_grad_critical.png}
  \caption{Critical points of the price slice $P(u_1, u_2 \mid u_3{=}0)$
    at $\gamma{=}100, \tau{=}1$: the zeros of $\nabla P$ (markers) lie
    on low-dimensional sets crossing the domain. Each critical value
    $p_c$ contributes an integrable singularity to the co-area weight
    and a cusp to the lookup.}
\end{figure}

\begin{figure}[H]\centering
  \includegraphics[width=0.8\linewidth]{cusps_figs/fig2_mu_cusps.png}
  \caption{The induced cusps in $\mu(p, u_k)$: spikes in the
    $p$-derivative align with the critical values of the slice map.
    Polynomial (global Chebyshev) representations ring around these
    cusps; shape-preserving (PCHIP) and \emph{segmented} representations
    do not.}
\end{figure}

\begin{figure}[H]\centering
  \includegraphics[width=0.7\linewidth]{cusps_figs/fig3_decay.png}
  \caption{Chebyshev coefficient decay of the equilibrium $P$ on Lobatto
    nodes at $\tau = 1$: algebraic ($\sim n^{-0.8}$), the quantitative
    signature of finite smoothness. Spectral methods cannot converge
    exponentially on this object without segmentation.}
\end{figure}

\chapter{Grids, coordinates, and precision}

\section{$u$-grid choices}
\begin{itemize}\itemsep2pt
  \item \textbf{CDF-uniform} ($u_i = \Phi^{-1}(q_i)$, mass-uniform):
        used by the kernel-band and \strict{} cube paths; right
        concentration for Gaussian-weighted integrals; \emph{wrong}
        nodes for polynomial fitting (Runge divergence at degree
        $\ge 13$, demonstrated).
  \item \textbf{Chebyshev--Lobatto}: correct interpolation nodes; used
        by the spectral experiments. The FP re-solved on Lobatto nodes
        converges to $|F| \sim 10^{-14}$ at $G = 11, 15, 21$, but the
        Chebyshev coefficients of the solution decay only
        algebraically (previous figure) --- the barrier is the
        function, not the nodes.
  \item \textbf{Private-log-odds uniform} ($\xi = \tanh(\tau u/2)$
        uniform, the ``no-learning posterior'' transform): the
        recommended representation grid (Part VI--VII). At
        $\xi_{\max} = 0.85$ the $u$-range is automatically
        $\approx \pm 5\sigma_{\rm signal}$ at every $\tau$ ---
        the grid \emph{scales itself} with signal informativeness, and
        no hard cube truncation is imposed.
\end{itemize}

\section{$p$-grid choices}
Logit-uniform is the workhorse. The adaptive arclength grid (Option D)
buys $\sim 25\%$ max-error at no cost. The \emph{public-log-odds} grid
($\Lpub(p) = \log[f_1(p)/f_0(p)]$ uniform, bootstrap-computed from the
cube CDF and symmetrized $\Lpub(1{-}p) = -\Lpub(p)$) is the natural
companion to the private-log-odds $u$-grid and is used by the final
configuration. Lesson learned: the bootstrap \emph{must} be symmetrized
--- numerical noise in independently computed $f_0, f_1$ densities broke
the sign-flip symmetry of the equilibrium by $O(0.4)$ until the
antisymmetrization $\Lpub \mapsto \tfrac12(\Lpub - \Lpub^{\rm rev})$ was
imposed by construction.

\section{Precision ladder}
Five arithmetic regimes were exercised across the project:
\begin{center}\begin{tabular}{lll}\toprule
regime & digits & role \\
\midrule
float64 & 16 & all production lookups; sufficient wherever floors are $\ge 10^{-12}$ \\
double-double (DD) & 32 & K=2 surface ladder; K=3 sweeps; the $2.7\times10^{-17}$-floor study \\
quad-double (QD) & 64 & spot checks \\
F128 & 34 & alternative to DD; $\sim 10\times$ slower, no accuracy gain --- retired \\
flint/mpmath & 100--130 & ground-truth anchors (\texttt{k3\_strict\_h0\_exact} nails toward $10^{-100}$) \\
\bottomrule\end{tabular}\end{center}
Working rule established by the DD-floor study: precision upgrades pay
only when the residual floor is at the arithmetic limit. The K=3
high-$\tau$ floors ($10^{-2}$) are operator-truncation floors, hence
\emph{not} improvable by DD/QD --- whereas the K=2 surface ladder and
the low-$\tau$ K=3 cells genuinely use the 32 digits.

\section{The K=2 fold at $\tau \approx 0.79$}
The K=2 double-double descent ($\tau: 2.0 \to 0.8$, $G = 32$) maps the
informative branch down to a fold: $\sigma_{\min}(I - J)$ collapses
$0.45 \to 0.31 \to 0.017$ across $\tau = 1.0, 0.9, 0.8$. Below the fold
the Newton descent cannot continue on this branch --- continuation in
$\tau$ alone is insufficient and pseudo-arclength would be required.
This is the earliest clean evidence in the project that the equilibrium
correspondence has folds, anticipating the K=3 branch-multiplicity
picture of Part VI.

\begin{figure}[H]\centering
  \includegraphics[width=0.85\linewidth]{fl_figs/fig1_F.png}
  \caption{The 1000-cell K=3 $\tau$-ladder residual under the
    kernel-band R4 operator (double-double arithmetic): machine-level
    floors at low $\tau$; the operator's intrinsic structure dominates
    above $\tau \sim 0.3$.}
\end{figure}

\begin{figure}[H]\centering
  \includegraphics[width=0.85\linewidth]{fl_figs/fig3_monotonicity.png}
  \caption{Monotonicity violations along the same ladder: the R4
    kernel-band fixed points violate axis-monotonicity in a large
    fraction of cells at high $\tau$ --- an operator artifact (the
    \strict{} and ansatz fixed points are clean), and an early warning
    that the R4 economics could not be taken at face value.}
\end{figure}

% =====================================================================
\part{The operator studies (Findings 1--4)}
% =====================================================================

\chapter{The kernel-band baseline and its failure (Findings 1--2)}

\section{The operator}
The session's starting production operator: Lin-CDF kernel bands at
$h \in \{0.5, 0.4, 0.3, 0.2\}$, combined by a 4-point Richardson
extrapolation in $h^2$ (``R4''), $\mathrm{NQK} = 16$ Gauss--Legendre
nodes per axis, evaluated at double-double fixed points on a $G = 11$
CDF-uniform grid with a $G_p = 121$ logit-uniform price grid.

\begin{figure}[H]\centering
  \includegraphics[width=0.9\linewidth]{ov_figs/fig1_richardson_ucurve.png}
  \caption{Richardson order U-curve at the saved R4 fixed point. R4 is
    locally optimal by construction; R3 under-cancels bias and R5
    over-amplifies noise (Vandermonde weight $L^1$-norm doubles per
    added $h$-point). A defensible choice \emph{conditional on the base
    method being right} --- which Finding 1 shows it is not.}
\end{figure}

\section{Finding 1: Richardson does not extrapolate to $h = 0$}
Bandwidths swept from 0.5 down to 0.02 with matched quadrature: the
lookup discrepancy versus \strict{} stays at 0.4--0.6 max-norm. If the
$h^2$ expansion were valid, the error would fall as $h^{2R}$; it does
not, because the level set of the piecewise-linear $P$ has kinks whose
contributions enter at non-integer powers $h^\alpha$ ($\alpha = 1$ or
$3/2$ depending on local geometry). Even-power Richardson cannot
annihilate them.

\begin{figure}[H]\centering
  \includegraphics[width=0.9\linewidth]{ov_figs/fig3_smallh.png}
  \caption{The smaller-$h$ test: the kernel-band lookup discrepancy vs
    \strict{} does not decay as $h \to 0.02$. The fixed-grid Richardson
    program is unfixable at its root.}
\end{figure}

\textbf{Reconciliation with Part II.} The kernel \emph{is} a consistent
co-area quadrature --- in the \emph{joint} limit $\Delta x \ll h \to 0$.
The failure here is the fixed-grid limit ($h \to 0$ at constant
$\Delta x$): once the band is thinner than the cell size, the
discretized integrand is kinky at the cell scale and the asymptotics in
$h$ break. The two statements are consistent and together they fix the
production rule: \emph{shrink $h$ only together with the grid}.

\section{Finding 2: the deficit collapse under \strict{}}
Re-solving the fixed point under \strict{} co-area + partition-of-unity,
warm-started from R4: the one-to-one deficit collapses by a factor
$\approx 40$ at every cell tested, e.g.\ at $(\gamma, \tau) = (100, 1)$
from $0.255$ to $0.0065$.

\begin{figure}[H]\centering
  \includegraphics[width=0.9\linewidth]{ov_figs/fig4_strict_vs_R4.png}
  \caption{R4 vs \strict{} fixed points at $(\gamma, \tau) = (100, 1)$:
    the R4 lookup systematically over-predicts the posterior at
    intermediate $p$.}
\end{figure}

\begin{figure}[H]\centering
  \includegraphics[width=0.9\linewidth]{ov_figs/fig5_mu_curves.png}
  \caption{The lookup at the \strict{} fixed point: smooth, monotone in
    $p$, mirror-symmetric about $p = \tfrac12$.}
\end{figure}

\section{The extended ridge recheck (36 cells)}
The K=3 ``deficit ridge'' (deficit non-monotone in $\tau$ at
intermediate $\gamma$) was re-examined on a $6 \times 6$ $(\gamma, \tau)$
grid under both operators. \textbf{The ridge does not survive}: under
\strict{} the deficit is flat in $\tau$ at $\sim 10^{-2}$, while R4
sweeps from $10^{-9}$ to $10^{-1}$. The operators \emph{cross} near
$\tau \approx 0.4$ --- R4 understates the deficit at low $\tau$
(undersmoothing of smooth cells) and overstates at high $\tau$
(oversmoothing of cuspy cells): a sign-flipping bias, not a uniform one.

\begin{figure}[H]\centering
  \includegraphics[width=\linewidth]{ridge_figs/slice_per_gamma.png}
  \caption{Deficit vs $\tau$ for $\gamma \in \{1, 10, 30\}$: R4 (red)
    rises steeply; \strict{} (blue) is flat. The crossing near
    $\tau \approx 0.4$ shows the bias flips sign.}
\end{figure}

\begin{figure}[H]\centering
  \includegraphics[width=0.8\linewidth]{ridge_figs/ratio_per_gamma.png}
  \caption{Collapse factor per cell: up to $\sim 30\times$ at high
    $\tau$, below 1 at low $\tau$.}
\end{figure}

\begin{figure}[H]\centering
  \includegraphics[width=\linewidth]{ridge_figs/per_cell_bars.png}
  \caption{Per-cell side-by-side deficits across the 18 completed
    low-$\gamma$ cells.}
\end{figure}

\textbf{Important caveat carried forward.} The \strict{} solver itself
is not Lipschitz (the cell-membership flips of Part II); Anderson
reaches $|F| \sim 0.03$--$0.15$ best iterates, not machine-eps fixed
points. The corrected deficit values are best-iterate quantities with
$\sim 10^{-3}$ noise floors. The truly nailed equilibria of the project
live in the consistent-kernel program (\texttt{k3\_coarea\_limit}, joint
central-path limit, $10^{-14}$), and the two views agree on the
qualitative geography: deficit rises with $\tau$, falls with $\gamma$.

\section{Finding 3: robustness $\ne$ correctness (the variant comparator)}
Seven perturbations of the kernel-band lookup --- PCHIP price
interpolation, Neville--Romberg, $\mathrm{NQK} = 24$, $G_p = 241$,
boundary refinements --- agree with each other to $\sim 10^{-5}$ while
\emph{all} sitting 0.05--0.6 from \strict{}. Cross-variant agreement
within an operator family measures conditioning, not truth.

\begin{figure}[H]\centering
  \includegraphics[width=0.95\linewidth]{ov_figs/fig2_variant_compare.png}
  \caption{The variant comparator: tight cross-variant spread, common
    bias. The defining picture of ``robust but wrong''.}
\end{figure}

\chapter{Alternative lookup builders (Finding 4)}

\section{The field}
Five algorithmically independent builders, evaluated against \strict{}
ground truth at matched fixed points:
\begin{center}\begin{tabularx}{\linewidth}{lXl}\toprule
option & construction & verdict \\
\midrule
A & 3D Chebyshev on CDF-uniform nodes + analytic level-set roots & Runge-divergent; wrong nodes \\
B+ & cube empirical CDF + PCHIP density ratio & \textbf{winner}: median $\sim 10^{-5}$ \\
C & structural ansatz $P = \sigma(\sum_i h(u_i))$ & exact at $\tau \to 0$; Part IV \\
D & adaptive $p$-grid (arclength equidistribution) & $-25\%$ max error, free \\
E & Hermite-cube derivative matching & no improvement \\
\bottomrule\end{tabularx}\end{center}

\begin{figure}[H]\centering
  \includegraphics[width=0.95\linewidth]{alts_figs/fig1_max_med.png}
  \caption{Max and median lookup error across builders and
    $(\gamma, \tau)$ cells: Option B+ improves the median by four
    decades over kernel-band R4.}
\end{figure}

\begin{figure}[H]\centering
  \includegraphics[width=0.95\linewidth]{alts_figs/fig2_mu_curves.png}
  \caption{Visual comparison of the builders' $\mu(p, u_k)$ against
    \strict{}: B+ tracks; A oscillates in the tails; C is exact at
    low $\tau$ but underfits at $\tau = 1$.}
\end{figure}

\section{Option B+ anatomy}
The construction is geometric, not extrapolative: accumulate
Gaussian-weighted cube-cell mass into an empirical CDF in $p$ per state,
differentiate with shape-preserving PCHIP, take the density ratio.
PCHIP's $C^1$-not-$C^\infty$ regularity is exactly matched to the
lookup's cusp structure --- it represents kinks without ringing.

\begin{figure}[H]\centering
  \includegraphics[width=0.95\linewidth]{alts_figs/fig3_cdf_construction.png}
  \caption{Option B+ construction: cube CDF, PCHIP density, resulting
    lookup vs \strict{}.}
\end{figure}

\begin{figure}[H]\centering
  \includegraphics[width=0.95\linewidth]{optB_figs/fig1_refinement.png}
  \caption{Sub-grid refinement: a $4\times$ refinement inside each bin
    buys roughly a decade of median accuracy at trivial cost (numba).}
\end{figure}

\begin{figure}[H]\centering
  \includegraphics[width=0.95\linewidth]{optB_figs/fig4_deficit_compare.png}
  \caption{The economic consequence: deficit maps under kernel-band R4
    vs Option B+ over the $(\gamma, \tau)$ region --- the spurious ridge
    versus the true pattern.}
\end{figure}

\textbf{B+ is an evaluator, not a solver.} Used as the inner operator of
a Picard loop it destabilizes the outer iteration --- it lacks the
contraction structure of the co-area operator. The division of labor
that survives to the final recommendation: co-area (kernel or strict)
\emph{solves}; B+/per-segment-Chebyshev \emph{evaluates}.

\section{Option D: the adaptive price grid}
Four monitor functions tested for $p$-grid equidistribution; arclength
of $\mu(p)$ wins (most noise-robust), reduces max interpolation error by
$\sim 25\%$ at $< 5\%$ cost, stacking multiplicatively with B+.

\begin{figure}[H]\centering
  \includegraphics[width=0.95\linewidth]{optD_figs/fig1_monitor_compare.png}
  \caption{Monitor-function comparison for the adaptive $p$-grid.}
\end{figure}

\begin{figure}[H]\centering
  \includegraphics[width=0.8\linewidth]{optD_figs/fig3_density.png}
  \caption{Resulting node-density profiles: the arclength monitor puts
    $\sim 23\times$ more nodes where the lookup actually bends.}
\end{figure}

\chapter{The spectral program and its limit}

\section{Chebyshev--Lobatto + analytic contours}
The elegant idea: fit a global 3D Chebyshev polynomial to $P$, then the
level set is an algebraic variety --- companion-matrix root-finding gives
exact contours and the lookup becomes clean 1D quadrature. Two results:
\begin{enumerate}\itemsep2pt
  \item On CDF-uniform nodes the fit diverges (Runge; edge coefficients
        $> 10^6$ at degree 15). Re-solving the FP on true Lobatto nodes
        fixes conditioning ($|F| \sim 10^{-14}$ at $G = 11, 15, 21$).
  \item But convergence of the \emph{lookup} is algebraic only:
        top-edge coefficient $\sim G^{-0.8}$, lookup median
        $\sim G^{-1.6}$. Machine epsilon requires effectively infinite
        $G$. The equilibrium is not analytic at high $\tau$ ---
        consistent with the $C^0$-not-$C^1$ co-area theorem of Part II.
\end{enumerate}

\begin{center}\begin{tabular}{rrrrr}\toprule
$G$ & FP $|F|_\infty$ & top edge coef & lookup max & lookup median \\
\midrule
11 & 6.8e-14 & 0.0491 & 0.434 & 0.251 \\
15 & 2.4e-14 & 0.0198 & 0.407 & 0.127 \\
21 & 4.3e-14 & 0.0110 & 0.380 & 0.081 \\
\bottomrule\end{tabular}\end{center}

\section{Ruling out the operator: the Chebyshev-native fixed point}
To exclude the possibility that the algebraic rate came from mixing the
spectral fit with the kernel-band operator, a Chebyshev-native FP solver
was built (lookup evaluated by analytic contours at the exact cell
price; no kernel, no Richardson, no bilinear interpolation). Same
algebraic decay. \textbf{The slow convergence is a property of the
equilibrium, not of any operator} --- the single most load-bearing
negative result of the spectral program, and the direct motivation for
the \emph{segmented} spectral representation that finally succeeds in
Part V.

% =====================================================================
\part{Structural ansatz and coordinate transformations (Findings 5--7)}
% =====================================================================

\chapter{Method C v2: the additive sigmoid ansatz}

\section{Motivation and form}
In the $\tau \to 0$ linear-Gaussian limit, signals enter posterior
log-odds additively and the (CARA/Hellwig) price is exactly
$\sigma(\tau \sum_i u_i)$. The structural ansatz preserves this skeleton
with a learned one-dimensional profile:
\[
  P(\mathbf{u}) = \sigma\bigl(h(u_1) + h(u_2) + h(u_3)\bigr),
  \qquad
  h(u) = \int_0^u s^2(t)\, dt,
  \qquad
  s(t) = \sum_{n=0}^{4} c_n\, T_{2n}(t / U_{\max}),
\]
five free coefficients. By construction: permutation-symmetric,
sign-flip symmetric, axis-monotone (zero violations across every run),
and equipped with a \emph{closed-form 1D lookup}: under the ansatz the
level set in $(u_a, u_b)$ at fixed $u_k$ collapses, after $v = h(u)$, to
the affine line $v_a + v_b = \operatorname{logit}p - h(u_k)$, making
$\mu(p, u_k)$ a single Gauss--Legendre integral
($\sim 10\,\mu$s per call --- four orders of magnitude faster than any
cube-based builder).

\begin{figure}[H]\centering
  \includegraphics[width=0.85\linewidth]{mcv2t_figs/fig1_sigma_h.png}
  \caption{The fitted profile at $\tau = 0.001$: $s^2(\cdot)$ and the
    resulting odd-monotone $h(u)$, matching the analytic
    linear-Gaussian limit $h(u) = \tau u + O(\tau^3 u^3)$ to plotting
    accuracy.}
\end{figure}

\section{The grids and ladders}
Trust-region least squares against \strict{} truth solves a cell in
$\sim 1$ second. Two production sweeps:
\begin{itemize}\itemsep2pt
  \item \textbf{100-cell $(\gamma, \tau)$ grid} (2 minutes): the
        residual is essentially a function of $\tau$ alone; at
        $\tau = 0.001$ the fit is exact to $10^{-10}$, at $\tau = 1$ it
        floors at $\|F\|_\infty \approx 0.08$--$0.10$.
  \item \textbf{1000-cell $\tau$-ladder} (16.5 minutes, chained
        warm-start, zero monotonicity violations): the residual scales
        as $\sim \tau^4$ at small $\tau$ (consistent with the
        $O(\tau^3 u^3)$ truncation of $h$ feeding an $O(\tau^4)$
        fixed-point error) and saturates at the non-additive plateau
        above $\tau \approx 0.3$--$0.4$.
\end{itemize}

\begin{figure}[H]\centering
  \includegraphics[width=0.9\linewidth]{mcv2_figs/fig1_F.png}
  \caption{Method C v2 1000-cell ladder: residual vs $\tau$. The sharp
    transition near $\tau \approx 0.3$ separates ``additive ansatz
    exact'' from ``non-additive correction required''.}
\end{figure}

\begin{figure}[H]\centering
  \includegraphics[width=0.9\linewidth]{mcv2_figs/fig2_loglog.png}
  \caption{Same data, log-log: the $\sim \tau^4$ scaling at small
    $\tau$ and the saturation plateau.}
\end{figure}

\begin{figure}[H]\centering
  \includegraphics[width=0.9\linewidth]{mcv2g_figs/fig1_F.png}
  \caption{The 100-cell $(\gamma, \tau)$ grid heat map: residual depends
    almost entirely on $\tau$, not $\gamma$ --- the additive structure
    is $\tau$-driven.}
\end{figure}

\section{Two control experiments}
\begin{itemize}\itemsep2pt
  \item \textbf{Logit-space loss}: refitting v2 with an
        $L$-space loss leaves the converged residual unchanged (median
        $1.7\times10^{-3}$ in both) --- the bottleneck is ansatz
        expressiveness, not loss geometry.
  \item \textbf{Strict-operator inner loop}: v2 fit with \strict{} as
        the inner operator at $\tau = 0.001$ converges to
        $P \equiv \tfrac12$ exactly --- the first clean sighting of the
        trivial branch that Part VI later confirms is an exact fixed
        point at all $(\gamma, \tau)$.
\end{itemize}

\chapter{Finding 5: the $(\Lpriv, \Lpub)$ coordinate test}

\section{Hypothesis}
Bayes suggests $\operatorname{logit}\mu(p, u_k) = \Lpriv(u_k) + \Lpub(p)
+ \delta(p, u_k)$ with $\Lpriv = \tau u_k$ (closed form),
$\Lpub = \log[f_1(p)/f_0(p)]$, and $\delta$ a small smooth K=3
conditional-overlap correction. If $\delta$ were small \emph{and}
smooth, Chebyshev in the transformed coordinates would converge
spectrally with a handful of coefficients.

\section{Result: bulk absorbed, cusps preserved}
The optimal additive (ANOVA) decomposition removes $\sim 50\%$ of the
RMS magnitude --- but the Chebyshev tail-decay of $\delta$ is
\emph{indistinguishable} from that of the raw lookup, and the SVD of
$\delta$ is as full-rank as the original. The cusps live on the
algebraic variety $\{|\nabla P| = 0\}$, which no global change of
coordinates can remove.

\begin{figure}[H]\centering
  \includegraphics[width=\linewidth]{lpub_figs/fig2_heat.png}
  \caption{$L$, $\delta_{\rm ANOVA}$, $\delta_{\rm Cheby-h}$ heat maps
    across four cells: the smooth vertical gradient is removed; the
    sharp horizontal cusp bands survive untouched.}
\end{figure}

\begin{figure}[H]\centering
  \includegraphics[width=\linewidth]{lpub_figs/fig1_tail_decay.png}
  \caption{Chebyshev tail-norm decay: the three curves (raw $L$, ANOVA
    residual, Cheby-$h$ residual) coincide. The coordinate
    transformation does not improve the convergence rate.}
\end{figure}

\begin{figure}[H]\centering
  \includegraphics[width=\linewidth]{lpub_figs/fig4_svd.png}
  \caption{SVD spectra: $\delta$ is essentially full rank; rank-1
    captures only 30--45\% of its Frobenius norm. No low-rank
    multiplicative correction exists either.}
\end{figure}

\section{The economic by-product}
The fitted $u_k$-effect $b(u_k)$ is roughly $10\times$ smaller than the
textbook $\tau u_k$: at $(\gamma, \tau) = (100, 1)$ the closed form
predicts $\pm 2.33$ where the data give $\pm 0.22$. Conditional on the
price, the agent's own signal adds little --- the K=3 price already does
most of the aggregation. This is the lookup-level reflection of the
Jensen-gap economics of Part I: the price is a powerful but
\emph{biased} (compressed) aggregator.

\begin{figure}[H]\centering
  \includegraphics[width=0.75\linewidth]{lpub_figs/fig3_effects.png}
  \caption{Row and column effects of the additive decomposition; dotted
    line: the naive $\tau u_k$ prediction, an order of magnitude too
    steep.}
\end{figure}

\chapter{Findings 6--7: the limits of finite-parameter structure}

\section{Finding 7: no low-degree polynomial closes the gap}
At $(\gamma, \tau) = (100, 1)$, the \emph{best achievable} additive fit
to the \strict{} log-odds surface leaves $|R|_\infty = 3.5$ (larger than
$|L|_\infty = 2.56$ itself), flat across 3--11 profile parameters; and
no symmetric polynomial correction captures the residual:

\begin{center}\begin{tabular}{lrr}\toprule
candidate $g$ & best $\kappa$ & share of $\|R\|^2$ captured \\
\midrule
$u_1 u_2 u_3$ & $-0.103$ & 7.8\% \\
$\sum_{i<j}(u_i - u_j)^2$ & $\approx 0$ & 0.0\% \\
$(u_1{+}u_2{+}u_3)^3$ & $-0.009$ & 22.8\% \\
$\sum u_i^3$, $\sum u_i^5$ & $\approx 0$ & 0.0\% \\
six-term LSQ combination & --- & 63.3\% \\
\bottomrule\end{tabular}\end{center}

\begin{figure}[H]\centering
  \includegraphics[width=0.95\linewidth]{mcv3_figs/refit.png}
  \caption{The fingerprint test: residual flat in profile flexibility
    (top left); the residual heat map (bottom right) is a high-rank
    checkerboard concentrated at corners --- not any smooth low-degree
    polynomial.}
\end{figure}

\section{Interpretation}
Findings 5 and 7 are the same fact seen from two angles: \textbf{the
non-additive structure of the K=3 equilibrium at high $\tau$ is
high-rank and cusp-bearing}. Neither clever coordinates nor
finite-parameter corrections reach $10^{-4}$; the viable routes are
(i) segmented representations that put the cusps on segment boundaries
(Part V --- this works), and (ii) brute resolution with a consistent
smooth operator (Part VII recommendation for the surface itself).

% =====================================================================
\part{The machine-epsilon lookup (Findings 8--16)}
% =====================================================================

\chapter{Cusp detection (Findings 8 and 13)}

\section{The 3D detector --- and why it is not enough}
A per-cube-cell critical-point finder (trilinear roots of
$\nabla \hat P = 0$, quadratic-Hermite fallback) enumerates the interior
critical points of the discretized price cube. At the four reference
\strict{} fixed points it finds 40--140 points (16--42 unique critical
values $p_c$), all with $|\nabla \hat P|$ at machine epsilon, symmetric
about $p = \tfrac12$, broadening with $\tau$.

\begin{figure}[H]\centering
  \includegraphics[width=0.95\linewidth]{critpts/dd_k3_strict_fp_g100_t1.0000.png}
  \caption{The 3D critical-point detector at
    $(\gamma, \tau) = (100, 1)$: 140 interior critical points; bimodal
    distribution of critical values; residuals at machine epsilon.}
\end{figure}

\textbf{Finding 13 (the correction that unlocked everything):} the
cusps of the \emph{one-dimensional} lookup $\mu(p, u_k)$ at fixed $u_k$
are critical values of the \emph{slice} map
$P(\cdot, \cdot, u_k)$ --- not of the full 3D map. Knotting segments at
the 3D critical values leaves a 3--9\% error floor at $\tau = 1$;
detecting the cusps directly from $\mu$ itself (second-difference
spikes plus first-difference jumps, pooled across $u_k$, deduplicated in
logit distance) finds the correct knots and collapses the floor to
machine epsilon. The lesson generalizes: \emph{locate singularities in
the object you are approximating, not in its generator}.

\section{Multi-element OptB+ (Finding 8): a half-step}
Inserting the (3D) critical values as forced PCHIP knots drives the
\emph{median} error to zero at the knots by construction but leaves the
max error unchanged --- because the max error was never at the bulk
cusps; it was in the tails ($p \to 0, 1$), where the empirical CDF runs
out of support. This diagnosis redirected the program to the tail
treatment of Findings 11--14.

\chapter{The segmented spectral representation (Finding 10)}

\section{Per-segment Chebyshev between cusp knots}
Between adjacent cusps the lookup is $C^\infty$; a degree-$d$ Chebyshev
fit per segment therefore converges exponentially \emph{per segment}.
With the correct (internal, $\mu$-based) knots:

\begin{center}\begin{tabular}{lrrr}\toprule
cell & degree & in-segment max & in-segment median \\
\midrule
$(100, 0.2)$ & 4 & 6.7e-16 & --- \\
$(100, 1.0)$ & 3 & 3.4e-15 & 3.3e-16 \\
$(1000, 0.2)$ & 4 & 6.7e-16 & 2.8e-16 \\
$(1000, 1.0)$ & 6 & 1.6e-15 & 2.2e-16 \\
\bottomrule\end{tabular}\end{center}

Machine epsilon at degree 3--6, all four reference cells --- the cusps
are honest piecewise-$C^k$ discontinuities at discrete $p_c$, not
smoothable curvature bumps. Notably, 67--92\% of segments contain
$\le 2$ samples (cusps cluster), forcing degree 0--1 on them; they reach
machine epsilon anyway because they are tiny in $p$-width.

\begin{figure}[H]\centering
  \includegraphics[width=0.75\linewidth]{pc_figs/converge_g100_t0.2000.png}
  \caption{Per-segment convergence at $(\gamma, \tau) = (100, 0.2)$:
    machine epsilon at any degree once knots are correct.}
\end{figure}

\begin{figure}[H]\centering
  \includegraphics[width=0.75\linewidth]{pc_figs/converge_g100_t1.0000.png}
  \caption{$(\gamma, \tau) = (100, 1.0)$: the in-segment median drops to
    $10^{-16}$ from degree 3; the blue max curve is dominated by
    out-of-support points (the fallback artifact addressed in
    Finding 14).}
\end{figure}

\begin{figure}[H]\centering
  \includegraphics[width=0.8\linewidth]{pc_figs/deg_heat_g100_t1.0000.png}
  \caption{Effective degree per $(u_k, \text{segment})$: large central
    segments support high degree; clustered cusp segments are forced
    low --- harmlessly.}
\end{figure}

\chapter{Negative control experiments (Findings 9, 11, 12, 17)}

\section{Finding 9: smoothing the strict operator has no sweet spot}
Replacing the discrete crossing test $ab < 0$ with the smooth weight
$\sigma(-ab/\varepsilon^2)$ makes the operator Lipschitz and Anderson
converges to machine epsilon for $\varepsilon \ge 0.1$ --- but the
resulting fixed point is biased 0.13--0.49 against the original
operator; and as $\varepsilon \to 0$ the convergence failure returns
before the bias disappears. There is no $\varepsilon$ at which both
residuals are small. The diagnostic that motivated this: \emph{starting
Picard at the saved ``strict fixed point'' diverges} (from
$|F| = 0.06$ to $0.5$) --- the saved object is a best iterate, not a
fixed point.

\begin{figure}[H]\centering
  \includegraphics[width=0.9\linewidth]{smooth_figs/eps_sweep.png}
  \caption{The $\varepsilon$-sweep: clean convergence under the smoothed
    operator (left) coexists with large bias under the original
    operator (right) at every $\varepsilon$.}
\end{figure}

\section{Findings 11--12: tail BCs and exact cell mass are not the levers}
Boundary-augmented PCHIP, logit-space CDF fitting, and Beta-CDF
parametric tails all leave the max error \emph{identical} to baseline
(the boundary error is interpolation failure, not missing BCs); exact
per-cell Gauss--Legendre mass integration halves the in-support median
but worsens the max (a smoother CDF extrapolates worse outside support).
Useful contributors, not regime changers.

\section{Finding 17: a smooth lookup does not fix the FP iteration}
The decisive control: using the machine-epsilon segmented lookup as the
\emph{inner operator} of the FP iteration, with cusp knots
\emph{frozen} (eliminating cusp drift, the conjectured source of
non-smoothness) still plateaus at $|F| \approx 1.1\times10^{-2}$ ---
essentially the same wall as the recompute mode (1.5e-2) and only
$5\times$ better than the kernel-band baseline (5.7e-2). The wall is the
$G = 11$ truncation error of the co-area evidence, not operator
roughness. PCHIP segments in the inner loop diverge outright.

\begin{figure}[H]\centering
  \includegraphics[width=\linewidth]{zhs_figs/fig1_convergence.png}
  \caption{The zero-h FP solver: recompute and frozen-knot modes plateau
    together at $\sim 10^{-2}$ while the kernel-band trace (grey)
    diverges. Smoothness of the lookup was not the binding constraint.}
\end{figure}

\begin{figure}[H]\centering
  \includegraphics[width=0.9\linewidth]{zhs_figs/fig2_mu_smoothing.png}
  \caption{Verification that the inner lookup is exact: the segmented
    pipeline reproduces all 369 in-support $\mu$ values at the
    warm-start to $2.2\times10^{-15}$.}
\end{figure}

\chapter{The production pipeline (Findings 14--16)}

\section{The fallback fix}
The raw \strict{} lookup returns $\mu = 0.5$ wherever the co-area
integral finds no roots ($p$ outside the cube's empirical price range),
producing $-0.5$ jumps and 0/11 monotone slices at high $\tau$. The fix
(\texttt{fix\_fallback}) replaces the fallback region by PCHIP
extrapolation through the in-support values pinned to the Bayesian
boundary conditions $\mu(\epsilon) = 0$, $\mu(1-\epsilon) = 1$:
afterwards, 11/11 slices monotone with correct tails at every cell
tested, and the spurious $|\mu''|$ spikes drop by half.

\begin{figure}[H]\centering
  \includegraphics[width=\linewidth]{v5_figs/g100_t1.0_deg6.png}
  \caption{The before/after of the production pipeline at
    $(\gamma, \tau) = (100, 1)$: raw $\mu_{\rm strict}$ with fallback
    discontinuities (left); cleaned (middle); pipeline v5 output ---
    monotone, smooth, boundary-correct (right).}
\end{figure}

\section{Pipeline v5 and the 36-cell certification}
v5 = \texttt{fix\_fallback} (clean) $\to$ internal cusp detector
(knots) $\to$ degree-6 Chebyshev per segment (in-support) $\to$
PCHIP shape-preserving tails (out-of-support). Applied to the full
$6 \times 6$ $(\gamma, \tau)$ ridge grid:

\begin{center}\begin{tabular}{lr}\toprule
metric & result \\
\midrule
cells processed & 36/36, 7 seconds total \\
median error vs cleaned truth & $1.11\times10^{-16}$ at every cell \\
max error at machine eps & 31/36 cells ($10^{-15}$--$10^{-14}$) \\
outlier cells & 5, max $10^{-2}$--$9\times10^{-2}$ (corner cells, fallback-smoothing artifacts) \\
monotone slices (dense 501-pt check) & 11/11 at every cell \\
boundary conditions & satisfied at every cell \\
\bottomrule\end{tabular}\end{center}

\begin{figure}[H]\centering
  \includegraphics[width=\linewidth]{v5g_figs/heatmap_v5.png}
  \caption{The certification heat maps: $\log_{10}$ max error (left),
    $\log_{10}$ median error --- uniformly $-16$ (middle), and monotone
    slice count --- uniformly 11/11 (right), across the
    $(\gamma, \tau)$ grid.}
\end{figure}

\section{What is and is not claimed}
The lookup problem --- evaluating the Bayes-correct conditional
posterior given a price surface --- is solved to machine epsilon
in-support, with controlled monotone Bayesian extrapolation
out-of-support, at one second per cell. The claim is conditional on the
input surface: it inherits whatever distance the underlying
(best-iterate) \strict{} surface has from the true equilibrium. The
surface question is Part VI's subject and drives the final
recommendation.

% =====================================================================
\part{The equilibrium surface: log-odds coordinates, branches, and solvers}
% =====================================================================

\chapter{The log-odds configuration}

\section{Specification}
The final coordinate system, run at the user's direction:
\begin{itemize}\itemsep2pt
  \item \textbf{$u$-grid}: uniform in the no-learning private posterior
        log-odds transform $\xi = \tanh(\tau u / 2)$,
        $\xi \in [-0.85, 0.85]$, mapped back by
        $u = (2/\tau)\operatorname{atanh}(\xi)$. No bounded cube: the
        grid self-scales to $\approx \pm 5\sigma_{\rm signal}$ at every
        $\tau$ (at $\tau = 0.005$ the cube spans $\pm 70$ in $u$; at
        $\tau = 1$, $\pm 2.4$).
  \item \textbf{$p$-grid}: uniform in public-only posterior log-odds
        $\Lpub(p)$, bootstrap-computed from the cube empirical CDF and
        \emph{antisymmetrized by construction}
        ($\Lpub \mapsto \tfrac12(\Lpub - \Lpub^{\rm rev})$) after the
        unsymmetrized bootstrap was caught breaking the equilibrium's
        sign-flip symmetry at $O(0.4)$.
  \item \textbf{Operator}: \strict{} co-area + POU (zero kernel).
  \item \textbf{Clearing}: exact CRRA bisection (the linearized
        clearing bug and its repair are recorded in Part I).
  \item \textbf{Residual}: both raw $|F|_\infty$ and the
        density-weighted $\Fw = \max |\Phi(P) - P| \cdot W_3$,
        $W_3 \propto \prod_i \bar f(u_i)$ --- the raw max over an
        unbounded cube is dominated by corner cells of vanishing
        probability and is not the economically meaningful residual.
  \item \textbf{Symmetry projection} (permutation + sign-flip) every
        iteration; live per-iteration reporting.
\end{itemize}

\section{The $G$-ladder and $\gamma$-sweep}
At $(\gamma, \tau) = (100, 0.1)$: $\Fw$ best of $1.3\times10^{-2}$
($G{=}7$) down to $1.6\times10^{-3}$ ($G{=}9$), non-monotone across
$G \in \{7..15\}$ --- the \strict{} stall manifests differently per
grid. The $\gamma$-sweep at $G = 11$ spans
$\Fw = 1.8\times10^{-2}$--$2.6\times10^{-2}$ across
$\gamma \in [1, 1000]$: genuinely $\gamma$-dependent (after the clearing
fix), all cells with sign-flip symmetry at $1.1\times10^{-16}$ and
$P(\mathbf{0}) = \tfrac12$ exactly.

\begin{figure}[H]\centering
  \includegraphics[width=0.9\linewidth]{lo_figs/g100_t0.1000_G11.png}
  \caption{The log-odds configuration at
    $(\gamma, \tau, G) = (100, 0.1, 11)$: raw vs density-weighted
    residual traces (left) and the equilibrium slice (right). The raw
    residual is dominated by zero-density corners; the weighted
    residual is the meaningful one.}
\end{figure}

\section{The 50-cell $\tau$-ladder}
Fifty log-spaced $\tau \in [0.001, 1.2]$, chained warm-start,
16 seconds total. Three regimes:
\[
  \underbrace{\tau \le 0.005}_{\Fw \approx 2.5\times10^{-15}}
  \quad\longrightarrow\quad
  \underbrace{0.005 \lesssim \tau \lesssim 0.02}_{\Fw: 10^{-7} \to 10^{-3}}
  \quad\longrightarrow\quad
  \underbrace{\tau \gtrsim 0.05}_{\Fw \approx 10^{-2} \text{ (stall)}}
\]
with perfect sign-flip symmetry everywhere and the price range opening
monotonically with $\tau$.

\begin{figure}[H]\centering
  \includegraphics[width=\linewidth]{lol_figs/ladder_summary.png}
  \caption{The 50-cell Picard $\tau$-ladder: residual regimes (left) and
    equilibrium price range vs $\tau$ (right).}
\end{figure}

\begin{figure}[H]\centering
  \includegraphics[width=\linewidth]{lol_figs/logodds_contours.png}
  \caption{Equilibrium log-odds contours
    $\operatorname{logit}P(u_1, u_2 \mid u_3{=}0)$ on private-log-odds
    axes across six $\tau$ values: a clean diagonal zero contour at low
    $\tau$ deforming into the kinked high-$\tau$ structure whose
    critical values generate the lookup cusps.}
\end{figure}

\begin{figure}[H]\centering
  \includegraphics[width=\linewidth]{lol_figs/logodds_contours_u.png}
  \caption{The same surfaces on raw $u$ axes with $k\sigma$ density
    rings: the unbounded-cube grid keeps the Gaussian mass interior at
    every $\tau$.}
\end{figure}

\section{Anatomy of one lookup evaluation}
For each $(p, u_k)$ query the \strict{} builder consumes: all $G^2$
slice vertex values (linear interpolation backbone), $2 \times 16$
Gauss--Legendre sweep nodes, and the level-set roots found per sweep
line --- typically crossing 25--56 of the $(G{-}1)^2$ cells. One full
table is $G_p \times G = 891$ such queries.

\begin{figure}[H]\centering
  \includegraphics[width=\linewidth]{lol_figs/lookup_points.png}
  \caption{Which grid points the lookup uses: cube vertices (black), GL
    sweep lines (light), level-set roots from each sweep (blue
    $\times$ / orange $+$), and the cells crossed by the level set
    (red). Top row: $\tau = 0.005$ --- the level set hugs the
    antidiagonal and off-median prices find no roots (the fallback
    region of Finding 14). Bottom: $\tau = 0.5$ --- well-resolved level
    sets across the cube.}
\end{figure}

\chapter{Branches and Newton-type solvers}

\section{The trivial branch is exact}
$P \equiv \tfrac12$ satisfies the K=3 CRRA clearing identically (all
posteriors equal the prior, all $R_i = 1$, zero demand): cold-started
solvers of every type converge to it instantly, to $F = 0$ exactly.
The equilibrium correspondence is multi-valued everywhere; \emph{branch
identification is part of the computation's burden of proof}. The
informative branch is reached only by continuation (warm-starting along
the $\tau$-ladder from the analytic low-$\tau$ solution).

\section{Newton--Krylov via the implicit function theorem}
The IFT Newton direction
$\delta P = -(\partial\Phi/\partial P - I)^{-1}(\Phi(P) - P)$,
implemented matrix-free by LGMRES with finite-difference Jacobian-vector
products and Wolfe line search, warm-started from the Picard ladder:
\begin{itemize}\itemsep2pt
  \item \textbf{Low $\tau$}: weighted residuals driven as deep as
        $10^{-25}$--$10^{-50}$ on cells where the basin is locally
        smooth --- far below the Picard floor.
  \item \textbf{High $\tau$}: plateaus at $\Fw \sim 10^{-3}$--$10^{-2}$
        with intermittent symmetry breaks --- the finite-difference
        Jacobian inherits the operator's $C^0$-not-$C^1$ roughness.
  \item \textbf{Cost}: 200--1400 $F$-evaluations per cell
        ($50$--$100\times$ Picard); total ladder 207\,s vs 16\,s.
\end{itemize}
MINPACK HYBRD (explicit numerical Jacobian + Broyden updates) was also
tried: in logit space it is captured by the trivial branch even from
informative warm starts; in $P$-space it exits the unit cube
(one cell reached $P \in [-1590, 1591]$) even with bound penalties.
Newton--Krylov with continuation warm starts is the only stable Newton
variant found.

\begin{figure}[H]\centering
  \includegraphics[width=\linewidth]{lon_figs/ladder_newton.png}
  \caption{The IFT-Newton ladder vs Picard: deeper convergence at low
    $\tau$, the same operator wall at high $\tau$, at $50$--$100\times$
    the cost. Newton refines; it does not rescue.}
\end{figure}

\section{Synthesis of the surface-side evidence}
Across Picard, Anderson, Newton--Krylov, HYBRD, smoothed-$\varepsilon$
operators, frozen-knot inner loops, and the structural ansatz: every
method that respects the \strict{} operator's discretization plateaus at
the same $G$-dependent wall at high $\tau$, and every method that
smooths past the wall pays in bias. Combined with the
$C^0$-not-$C^1$ theorem and the success of the consistent-kernel joint
limit at $10^{-14}$ (\texttt{k3\_coarea\_limit}), the conclusion is
structural: \textbf{the high-resolution surface computation must use a
consistent $C^\infty$ operator (kernel co-area in the central-path
limit), with the bare \strict{} operator reserved for the role of
unbiased cross-check at matched tolerance}. That division is the spine
of the recommendation in Part VII.

% =====================================================================
\part{The recommendation: a reproducible \strict{}-anchored gold standard}
% =====================================================================

\chapter{Design principles}

The recommendation reconciles two facts that pull in opposite
directions:
\begin{enumerate}\itemsep2pt
  \item The \emph{bare} \strict{} operator is the unbiased definition
        of the equilibrium (no auxiliary parameter at all) --- but it is
        $C^0$-not-$C^1$ at Morse-critical prices and therefore cannot
        be nailed at high resolution by any Newton-type method
        (Part II theorem; Findings 9, 17; Part VI plateaus).
  \item The \emph{consistent kernel} co-area operator is $C^\infty$ and
        nails the equilibrium to $10^{-14}$ --- but carries the
        bandwidth $h$, which must be demonstrated (not assumed) to be a
        pure convergence knob.
\end{enumerate}
The resolution: \textbf{solve with the kernel in the joint central-path
limit; certify against \strict{} at matched tolerance; report the
$h$-ladder as part of the result}. The kernel is then never trusted ---
it is verified, cell by cell.

\chapter{The protocol, step by step}

\section{Step 0: coordinates and grids}
\begin{itemize}\itemsep2pt
  \item $u$-grid: uniform private log-odds $\xi = \tanh(\tau u/2)$,
        $\xi_{\max} = 0.85$ (auto-scaling $\pm 5\sigma$; no hard
        truncation). $G$ per the accuracy target below.
  \item $p$-grid: uniform public log-odds $\Lpub$, bootstrap from the
        cube CDF, antisymmetrized by construction; refreshed during
        warm-up only (iteration 5), frozen afterwards.
  \item Weights: trapezoid in $u$; density weight
        $W_3 = \prod_i \bar f(u_i)$ cached for the residual metric.
\end{itemize}

\section{Step 1: continuation ladder (the branch certificate)}
Start at $\tau_0 = 10^{-3}$ from the analytic additive solution
$P = \sigma(\tau \sum_i u_i)$ (Method C v2 closed form; exact there to
$10^{-10}$). March $\tau$ upward in log steps of $\le 5\%$, each cell
warm-started from its predecessor, recording per cell: $\Fw$ best,
symmetry errors, price range, wall time. This ladder \emph{is} the
proof of branch identity: a connected continuation path from the
analytic limit to the target cell. Never accept a cold-started
high-$\tau$ fixed point --- the trivial branch is exact and
attracting.

\section{Step 2: solve each cell with the consistent kernel, joint limit}
At each $(\gamma, \tau)$ target:
\begin{enumerate}\itemsep2pt
  \item Solve at $(G_1, h_1)$ with the Gaussian-band co-area operator,
        Anderson or Picard, to its floor; require the floor to be at
        arithmetic level for the precision in use.
  \item Refine along the central path: $(G_{k+1}, h_{k+1}) =
        (\lceil 1.4 G_k \rceil,\; h_k / 1.4)$ keeping
        $h \gtrsim 3\Delta\xi$ always, warm-starting each rung by
        interpolation. Three rungs minimum.
  \item Convergence claim: economic quantities (deficit, slope,
        $P$-range) must agree across the last two rungs to the target
        tolerance, demonstrating $h$-invariance. \emph{The ladder
        table, not a single number, is the result.}
  \item Symmetry projection (permutation + sign-flip) every iteration;
        symmetry errors reported at arithmetic level.
  \item Precision: float64 until the rung floor reaches $10^{-13}$;
        double-double beyond. (QD/F128 retired: no benefit at these
        floors.)
\end{enumerate}

\section{Step 3: the \strict{} cross-check (the unbiasedness certificate)}
At the finest rung, evaluate the \emph{bare} \strict{} residual at the
kernel solution. Accept if
$|F_{\rm strict}|_w \le C \cdot (\text{kernel-rung truncation
estimate})$ with $C \approx 3$. This certifies that the kernel solution
is a \strict{} near-fixed-point to the accuracy the bare operator can
measure --- the kernel never substitutes its own notion of equilibrium.

\section{Step 4: extract the lookup with pipeline v5}
On the certified surface: \texttt{fix\_fallback} $\to$ internal
$\mu$-based cusp detector $\to$ degree-6 per-segment Chebyshev $\to$
PCHIP monotone tails with $\mu(\epsilon) = 0$, $\mu(1-\epsilon) = 1$.
Certify per cell: in-support max error vs the cleaned co-area lookup at
machine epsilon; 11/11 monotone slices on a dense grid; boundary
conditions met. (Demonstrated at 36/36 cells, 7 seconds total.)

\section{Step 5: the validation battery (all five, every cell)}
\begin{enumerate}\itemsep2pt
  \item \textbf{Symmetry}: permutation, sign-flip, lookup mirror at
        arithmetic level.
  \item \textbf{Monotonicity}: $P$ axis-monotone; $\mu$ monotone in $p$
        per slice (dense check).
  \item \textbf{Boundary conditions}: lookup tails; $P(\mathbf{0}) =
        \tfrac12$ exactly.
  \item \textbf{Operator cross-check}: Step 3 inequality, plus the
        kernel $h$-ladder invariance table.
  \item \textbf{Measure robustness}: report the deficit under at least
        two state weightings (uniform and signal-weighted); the
        $\sim 0.12$ measure spread is part of the honest result, the
        price \emph{surface} being the invariant object.
\end{enumerate}

\section{Step 6: reporting standard}
Every published cell carries: the continuation provenance (ladder
path), the final rung $(G, h, \text{precision})$, the kernel-ladder
invariance table, the \strict{} cross-check value, best-iterate vs
converged status, the symmetry residuals, and the deficit with its
measure convention. Nothing else counts as a result.

\chapter{Accuracy targets and cost}

\begin{center}\begin{tabularx}{\linewidth}{lXXX}\toprule
target & grid/rungs & precision & cost/cell (order) \\
\midrule
survey ($10^{-3}$ econ) & $G{=}11$, 1 rung + v5 lookup & float64 & seconds \\
publication ($10^{-6}$ econ) & $G{=}17{-}25$, 3 rungs & float64 & minutes \\
anchor ($10^{-12+}$) & $G{\ge}31$, 4 rungs & DD (\& flint spot checks) & hours \\
\bottomrule\end{tabularx}\end{center}

The existing anchors: \texttt{k3\_coarea\_limit} ($10^{-14}$, joint
limit), \texttt{k3\_strict\_h0\_exact} (flint, toward $10^{-100}$ at
coarse $G$), and the $G{=}9$ no-kernel nail at $8\times10^{-15}$ ---
three independent operator families agreeing on the same PR equilibrium
is the strongest truth-closeness evidence the project owns.

\chapter{Open problems}

\begin{enumerate}\itemsep2pt
  \item \textbf{Analytic Jacobian through the kernel co-area operator}
        ($\sim 200$ LOC chain rule): would replace finite-difference
        Newton noise and likely extend deep convergence to higher
        $\tau$ at fixed cost.
  \item \textbf{Pseudo-arclength continuation} past the K=2 fold at
        $\tau \approx 0.79$ and a systematic K=3 fold map: the branch
        structure above $\tau \sim 1$ is mapped only along ladders, not
        in the $(\gamma, \tau)$ plane.
  \item \textbf{The 5 corner outlier cells} of the v5 certification
        ($10^{-2}$ max): localized fallback-smoothing artifacts at
        extreme parameter corners; a targeted $G_p$ refinement there is
        unfinished.
  \item \textbf{K=4 extreme-heterogeneity PR} (\texttt{k4\_pr\_nail}
        floors at $2.6\times10^{-3}$): whether it nails fully under the
        consistent co-area operator is open --- the natural next
        economics result.
  \item \textbf{A cusp-tracking surface representation}: the lookup-side
        segmentation idea lifted to the surface itself (multi-element
        cube patches bounded by the critical manifold) --- the
        principled route to exponential convergence of $P$, beyond the
        scope of the present cycle.
\end{enumerate}


% =====================================================================
\appendix
\part*{Appendices}
\addcontentsline{toc}{part}{Appendices}
\chapter{Per-cell result tables}

\section{The 50-cell log-odds Picard ladder ($\gamma=100$, $G=11$)}
{\small\begin{longtable}{rrrrrr}
\toprule
idx & $\tau$ & $\Fw$ best & $P_{\min}$ & $P_{\max}$ & sym err \\
\midrule
\endhead
0 & 0.00100 & 2.55e-15 & 0.0005 & 0.9995 & 4.4e-16 \\
1 & 0.00116 & 2.55e-15 & 0.0005 & 0.9995 & 4.4e-16 \\
2 & 0.00134 & 2.55e-15 & 0.0005 & 0.9995 & 4.4e-16 \\
3 & 0.00154 & 2.55e-15 & 0.0005 & 0.9995 & 4.4e-16 \\
4 & 0.00178 & 2.55e-15 & 0.4347 & 0.5653 & 1.1e-16 \\
5 & 0.00206 & 1.38e-14 & 0.3832 & 0.6168 & 1.1e-16 \\
6 & 0.00238 & 4.08e-13 & 0.3693 & 0.6307 & 1.1e-16 \\
7 & 0.00275 & 5.69e-07 & 0.3491 & 0.6509 & 1.1e-16 \\
8 & 0.00318 & 1.76e-02 & 0.4162 & 0.5838 & 1.1e-16 \\
9 & 0.00368 & 7.37e-07 & 0.4162 & 0.5838 & 1.1e-16 \\
10 & 0.00425 & 4.91e-08 & 0.4270 & 0.5730 & 1.1e-16 \\
11 & 0.00491 & 5.57e-08 & 0.3346 & 0.6654 & 1.1e-16 \\
12 & 0.00568 & 3.77e-07 & 0.3346 & 0.6654 & 1.1e-16 \\
13 & 0.00656 & 1.27e-06 & 0.3833 & 0.6167 & 1.1e-16 \\
14 & 0.00758 & 1.98e-02 & 0.3593 & 0.6407 & 1.1e-16 \\
15 & 0.00876 & 7.89e-05 & 0.3646 & 0.6354 & 1.1e-16 \\
16 & 0.01013 & 1.71e-04 & 0.3925 & 0.6075 & 1.1e-16 \\
17 & 0.01170 & 2.86e-04 & 0.3552 & 0.6448 & 1.1e-16 \\
18 & 0.01352 & 5.53e-04 & 0.3543 & 0.6457 & 1.1e-16 \\
19 & 0.01563 & 6.53e-04 & 0.2897 & 0.7103 & 1.1e-16 \\
20 & 0.01806 & 1.39e-03 & 0.3298 & 0.6702 & 1.1e-16 \\
21 & 0.02088 & 2.41e-04 & 0.2741 & 0.7259 & 1.1e-16 \\
22 & 0.02413 & 5.53e-02 & 0.2741 & 0.7259 & 1.1e-16 \\
23 & 0.02788 & 4.91e-03 & 0.3139 & 0.6861 & 1.1e-16 \\
24 & 0.03222 & 6.53e-03 & 0.3139 & 0.6861 & 1.1e-16 \\
25 & 0.03724 & 8.35e-03 & 0.3139 & 0.6861 & 1.1e-16 \\
26 & 0.04304 & 1.06e-02 & 0.3139 & 0.6861 & 1.1e-16 \\
27 & 0.04974 & 1.22e-02 & 0.3276 & 0.6724 & 1.1e-16 \\
28 & 0.05748 & 1.52e-02 & 0.1599 & 0.8401 & 1.1e-16 \\
29 & 0.06643 & 1.66e-02 & 0.2509 & 0.7491 & 1.1e-16 \\
30 & 0.07677 & 1.98e-02 & 0.2370 & 0.7630 & 1.1e-16 \\
31 & 0.08873 & 1.86e-02 & 0.2205 & 0.7795 & 1.1e-16 \\
32 & 0.10254 & 2.13e-02 & 0.2491 & 0.7509 & 1.1e-16 \\
33 & 0.11850 & 1.58e-02 & 0.2083 & 0.7917 & 1.1e-16 \\
34 & 0.13695 & 1.74e-02 & 0.2244 & 0.7756 & 1.1e-16 \\
35 & 0.15828 & 2.04e-02 & 0.2244 & 0.7756 & 1.1e-16 \\
36 & 0.18292 & 2.31e-02 & 0.2244 & 0.7756 & 1.1e-16 \\
37 & 0.21140 & 2.59e-02 & 0.1977 & 0.8023 & 1.1e-16 \\
38 & 0.24431 & 3.31e-02 & 0.2207 & 0.7793 & 1.1e-16 \\
39 & 0.28234 & 3.58e-02 & 0.2159 & 0.7841 & 1.1e-16 \\
40 & 0.32630 & 3.71e-02 & 0.2362 & 0.7638 & 1.1e-16 \\
41 & 0.37710 & 4.28e-02 & 0.2280 & 0.7720 & 1.1e-16 \\
42 & 0.43581 & 4.54e-02 & 0.2046 & 0.7954 & 1.1e-16 \\
43 & 0.50366 & 5.95e-02 & 0.2587 & 0.7413 & 1.1e-16 \\
44 & 0.58208 & 6.19e-02 & 0.2271 & 0.7729 & 1.1e-16 \\
45 & 0.67270 & 6.87e-02 & 0.2271 & 0.7729 & 1.1e-16 \\
46 & 0.77743 & 2.75e-02 & 0.3324 & 0.6676 & 1.1e-16 \\
47 & 0.89846 & 9.99e-02 & 0.2471 & 0.7529 & 1.1e-16 \\
48 & 1.03834 & 1.26e-01 & 0.2121 & 0.7879 & 1.1e-16 \\
49 & 1.20000 & 1.56e-01 & 0.2308 & 0.7692 & 1.1e-16 \\
\bottomrule
\end{longtable}}

\section{The 50-cell IFT-Newton ladder (warm-started from Picard)}
{\small\begin{longtable}{rrrrrrr}
\toprule
idx & $\tau$ & $F_\infty$ & $\Fw$ & $P_{\min}$ & $P_{\max}$ & nfev \\
\midrule
\endhead
0 & 0.00100 & 3.33e-16 & 1.46e-67 & 0.5000 & 0.5000 & 27 \\
1 & 0.00116 & 3.33e-16 & 1.15e-60 & 0.5000 & 0.5000 & 27 \\
2 & 0.00134 & 3.33e-16 & 1.07e-54 & 0.5000 & 0.5000 & 27 \\
3 & 0.00154 & 3.33e-16 & 1.56e-49 & 0.5000 & 0.5000 & 27 \\
4 & 0.00178 & 9.59e-13 & 2.02e-29 & 0.4398 & 0.5602 & 156 \\
5 & 0.00206 & 4.80e-02 & 6.31e-18 & 0.3855 & 0.6145 & 837 \\
6 & 0.00238 & 6.22e-02 & 2.37e-13 & 0.3838 & 0.6231 & 297 \\
7 & 0.00275 & 3.22e-15 & 3.13e-25 & 0.3735 & 0.6265 & 468 \\
8 & 0.00318 & 2.84e-01 & 2.66e-10 & 0.3731 & 0.6345 & 507 \\
9 & 0.00368 & 1.54e-01 & 3.15e-11 & 0.3111 & 0.6889 & 823 \\
10 & 0.00425 & 2.02e-01 & 6.87e-10 & 0.2330 & 0.7668 & 639 \\
11 & 0.00491 & 1.27e-01 & 1.12e-07 & 0.3425 & 0.6578 & 587 \\
12 & 0.00568 & 2.70e-02 & 8.30e-08 & 0.3205 & 0.6795 & 609 \\
13 & 0.00656 & 1.61e-01 & 5.03e-07 & 0.3064 & 0.7042 & 598 \\
14 & 0.00758 & 1.14e-02 & 4.05e-09 & 0.3733 & 0.6267 & 761 \\
15 & 0.00876 & 4.56e-02 & 1.50e-06 & 0.3648 & 0.6352 & 548 \\
16 & 0.01013 & 1.41e-01 & 4.97e-05 & 0.2567 & 0.7433 & 548 \\
17 & 0.01170 & 2.15e-02 & 3.57e-06 & 0.3578 & 0.6422 & 998 \\
18 & 0.01352 & 6.86e-02 & 1.24e-04 & 0.3708 & 0.6292 & 861 \\
19 & 0.01563 & 3.87e-02 & 7.32e-05 & 0.3361 & 0.6639 & 626 \\
20 & 0.01806 & 1.19e-01 & 4.97e-05 & 0.2859 & 0.7141 & 794 \\
21 & 0.02088 & 3.61e-02 & 1.31e-04 & 0.3847 & 0.6153 & 857 \\
22 & 0.02413 & 5.69e-02 & 3.10e-04 & 0.3866 & 0.6134 & 654 \\
23 & 0.02788 & 3.33e-02 & 1.34e-03 & 0.3202 & 0.6798 & 697 \\
24 & 0.03222 & 1.13e-01 & 1.21e-03 & 0.3534 & 0.6464 & 651 \\
25 & 0.03724 & 1.38e-01 & 6.84e-03 & 0.2757 & 0.7235 & 528 \\
26 & 0.04304 & 1.86e-01 & 1.67e-03 & 0.3122 & 0.6878 & 554 \\
27 & 0.04974 & 9.35e-02 & 2.74e-03 & 0.2818 & 0.7180 & 841 \\
28 & 0.05748 & 1.40e-01 & 1.93e-03 & 0.2839 & 0.7157 & 582 \\
29 & 0.06643 & 7.64e-02 & 7.74e-04 & 0.2387 & 0.7614 & 927 \\
30 & 0.07677 & 1.71e-01 & 4.89e-03 & 0.1604 & 0.8395 & 623 \\
31 & 0.08873 & 4.54e-02 & 1.28e-03 & 0.2056 & 0.7944 & 939 \\
32 & 0.10254 & 4.57e-02 & 8.03e-04 & 0.2418 & 0.7582 & 1392 \\
33 & 0.11850 & 2.93e-02 & 9.77e-04 & 0.2481 & 0.7519 & 1369 \\
34 & 0.13695 & 2.83e-02 & 2.52e-03 & 0.2303 & 0.7697 & 1262 \\
35 & 0.15828 & 1.45e-01 & 2.75e-03 & 0.2142 & 0.7858 & 878 \\
36 & 0.18292 & 1.26e-01 & 1.79e-02 & 0.1800 & 0.8200 & 1232 \\
37 & 0.21140 & 8.21e-02 & 8.63e-03 & 0.1943 & 0.8057 & 1127 \\
38 & 0.24431 & 9.06e-02 & 4.80e-03 & 0.2119 & 0.7881 & 959 \\
39 & 0.28234 & 1.17e-01 & 3.95e-03 & 0.2210 & 0.7790 & 1098 \\
40 & 0.32630 & 2.83e-02 & 2.35e-03 & 0.2246 & 0.7754 & 1274 \\
41 & 0.37710 & 8.25e-02 & 2.65e-02 & 0.1959 & 0.8041 & 1338 \\
42 & 0.43581 & 7.24e-02 & 8.82e-03 & 0.2587 & 0.7418 & 851 \\
43 & 0.50366 & 6.59e-02 & 1.12e-02 & 0.2430 & 0.7570 & 1431 \\
44 & 0.58208 & 2.13e-01 & 3.94e-02 & 0.1786 & 0.8214 & 390 \\
45 & 0.67270 & 7.29e-02 & 8.38e-03 & 0.2052 & 0.7950 & 940 \\
46 & 0.77743 & 2.01e-01 & 8.57e-02 & 0.1936 & 0.8046 & 664 \\
47 & 0.89846 & 2.56e-01 & 8.64e-02 & 0.0782 & 0.7935 & 592 \\
48 & 1.03834 & 1.03e-02 & 5.04e-03 & 0.2586 & 0.7414 & 1282 \\
49 & 1.20000 & 2.09e-01 & 1.02e-02 & 0.2252 & 0.7748 & 997 \\
\bottomrule
\end{longtable}}

\section{Pipeline v5 certification, 36 ridge cells}
{\small\begin{longtable}{rrrrrrr}
\toprule
$\gamma$ & $\tau$ & cusps & max err & median err & monotone & min diff \\
\midrule
\endhead
1 & 0.1 & 83 & 1.55e-15 & 1.11e-16 & 11/11 & 5.4e-04 \\
1 & 0.2 & 119 & 5.55e-16 & 1.11e-16 & 11/11 & 0.0e+00 \\
1 & 0.3 & 56 & 2.11e-15 & 1.11e-16 & 11/11 & 0.0e+00 \\
1 & 0.5 & 44 & 2.11e-15 & 1.11e-16 & 11/11 & 0.0e+00 \\
1 & 0.7 & 117 & 9.37e-02 & 1.11e-16 & 11/11 & -1.1e-16 \\
1 & 1.0 & 45 & 4.82e-02 & 1.11e-16 & 11/11 & -1.4e-17 \\
10 & 0.1 & 88 & 2.11e-15 & 1.11e-16 & 11/11 & 7.7e-04 \\
10 & 0.2 & 55 & 3.72e-02 & 1.11e-16 & 11/11 & 1.2e-05 \\
10 & 0.3 & 55 & 2.55e-15 & 1.11e-16 & 11/11 & 1.5e-05 \\
10 & 0.5 & 41 & 1.22e-15 & 1.11e-16 & 11/11 & 0.0e+00 \\
10 & 0.7 & 51 & 3.22e-15 & 1.11e-16 & 11/11 & -1.1e-16 \\
10 & 1.0 & 62 & 2.22e-15 & 1.11e-16 & 11/11 & -2.2e-16 \\
30 & 0.1 & 52 & 1.55e-15 & 1.11e-16 & 11/11 & 1.7e-03 \\
30 & 0.2 & 67 & 3.84e-02 & 1.11e-16 & 11/11 & 1.2e-05 \\
30 & 0.3 & 55 & 8.88e-16 & 1.11e-16 & 11/11 & 0.0e+00 \\
30 & 0.5 & 52 & 9.99e-16 & 1.11e-16 & 11/11 & 0.0e+00 \\
30 & 0.7 & 44 & 2.55e-15 & 1.11e-16 & 11/11 & -1.1e-16 \\
30 & 1.0 & 54 & 1.78e-15 & 1.11e-16 & 11/11 & 0.0e+00 \\
100 & 0.1 & 49 & 1.33e-15 & 1.11e-16 & 11/11 & 7.4e-04 \\
100 & 0.2 & 47 & 2.66e-15 & 1.11e-16 & 11/11 & 1.2e-05 \\
100 & 0.3 & 60 & 2.11e-15 & 1.11e-16 & 11/11 & 9.0e-06 \\
100 & 0.5 & 51 & 2.44e-15 & 1.11e-16 & 11/11 & 0.0e+00 \\
100 & 0.7 & 42 & 2.44e-15 & 1.11e-16 & 11/11 & -1.1e-16 \\
100 & 1.0 & 61 & 3.22e-02 & 1.11e-16 & 11/11 & 0.0e+00 \\
300 & 0.1 & 119 & 5.55e-16 & 1.11e-16 & 11/11 & 5.0e-04 \\
300 & 0.2 & 79 & 5.55e-16 & 1.11e-16 & 11/11 & 1.4e-05 \\
300 & 0.3 & 70 & 2.00e-15 & 1.11e-16 & 11/11 & 0.0e+00 \\
300 & 0.5 & 39 & 2.33e-15 & 1.11e-16 & 11/11 & 0.0e+00 \\
300 & 0.7 & 47 & 1.67e-15 & 1.11e-16 & 11/11 & -1.1e-16 \\
300 & 1.0 & 50 & 3.17e-04 & 1.11e-16 & 11/11 & 0.0e+00 \\
1000 & 0.1 & 52 & 1.11e-15 & 1.11e-16 & 11/11 & 9.7e-04 \\
1000 & 0.2 & 53 & 1.11e-15 & 1.11e-16 & 11/11 & 1.2e-05 \\
1000 & 0.3 & 53 & 4.11e-15 & 1.11e-16 & 11/11 & 0.0e+00 \\
1000 & 0.5 & 45 & 1.55e-15 & 1.11e-16 & 11/11 & 5.5e-06 \\
1000 & 0.7 & 46 & 1.78e-15 & 1.11e-16 & 11/11 & 0.0e+00 \\
1000 & 1.0 & 52 & 8.88e-16 & 1.11e-16 & 11/11 & -1.1e-16 \\
\bottomrule
\end{longtable}}

\section{Method C v2: the 100-cell $(\gamma,\tau)$ grid}
{\small\begin{longtable}{rrrrrr}
\toprule
$\gamma$ & $\tau$ & $\|F\|_\infty$ & slope & $P_{\min}$ & $P_{\max}$ \\
\midrule
\endhead
0.1 & 0.001 & 1.24e-08 & 0.0003 & 0.4994 & 0.5006 \\
0.1 & 0.005 & 1.72e-06 & 0.0017 & 0.4971 & 0.5029 \\
0.1 & 0.01 & 1.37e-05 & 0.0033 & 0.4942 & 0.5058 \\
0.1 & 0.05 & 1.64e-03 & 0.0167 & 0.4710 & 0.5290 \\
0.1 & 0.1 & 1.05e-02 & 0.0337 & 0.4424 & 0.5576 \\
0.1 & 0.3 & 5.91e-02 & 0.1123 & 0.3241 & 0.6759 \\
0.1 & 0.5 & 4.16e-02 & 0.2406 & 0.1812 & 0.8188 \\
0.1 & 0.7 & 3.27e-02 & 0.3627 & 0.0966 & 0.9034 \\
0.1 & 0.85 & 5.43e-02 & 0.6891 & 0.0075 & 0.9925 \\
0.1 & 1 & 7.90e-02 & 0.9228 & 0.0010 & 0.9990 \\
0.3 & 0.001 & 2.73e-09 & 0.0003 & 0.4994 & 0.5006 \\
0.3 & 0.005 & 1.41e-07 & 0.0017 & 0.4971 & 0.5029 \\
0.3 & 0.01 & 1.14e-06 & 0.0033 & 0.4942 & 0.5058 \\
0.3 & 0.05 & 1.39e-04 & 0.0167 & 0.4709 & 0.5291 \\
0.3 & 0.1 & 1.04e-03 & 0.0338 & 0.4414 & 0.5586 \\
0.3 & 0.3 & 1.28e-02 & 0.1149 & 0.3135 & 0.6865 \\
0.3 & 0.5 & 3.32e-03 & 0.2676 & 0.1418 & 0.8582 \\
0.3 & 0.7 & 4.55e-02 & 0.5272 & 0.0246 & 0.9754 \\
0.3 & 0.85 & 5.81e-02 & 0.7410 & 0.0052 & 0.9948 \\
0.3 & 1 & 8.83e-02 & 0.9479 & 0.0011 & 0.9989 \\
1 & 0.001 & 1.78e-10 & 0.0003 & 0.4994 & 0.5006 \\
1 & 0.005 & 1.88e-08 & 0.0017 & 0.4971 & 0.5029 \\
1 & 0.01 & 1.72e-07 & 0.0033 & 0.4942 & 0.5058 \\
1 & 0.05 & 2.14e-05 & 0.0168 & 0.4708 & 0.5292 \\
1 & 0.1 & 1.64e-04 & 0.0340 & 0.4410 & 0.5590 \\
1 & 0.3 & 1.76e-03 & 0.1200 & 0.3029 & 0.6971 \\
1 & 0.5 & 2.47e-02 & 0.2989 & 0.1116 & 0.8884 \\
1 & 0.7 & 5.68e-02 & 0.5738 & 0.0169 & 0.9831 \\
1 & 0.85 & 5.89e-02 & 0.7966 & 0.0036 & 0.9964 \\
1 & 1 & 1.02e-01 & 1.0132 & 0.0007 & 0.9993 \\
3 & 0.001 & 2.33e-09 & 0.0003 & 0.4994 & 0.5006 \\
3 & 0.005 & 5.94e-09 & 0.0017 & 0.4971 & 0.5029 \\
3 & 0.01 & 5.76e-08 & 0.0033 & 0.4942 & 0.5058 \\
3 & 0.05 & 7.02e-06 & 0.0168 & 0.4708 & 0.5292 \\
3 & 0.1 & 5.01e-05 & 0.0341 & 0.4409 & 0.5591 \\
3 & 0.3 & 1.12e-03 & 0.1220 & 0.2994 & 0.7006 \\
3 & 0.5 & 3.12e-02 & 0.3086 & 0.1039 & 0.8961 \\
3 & 0.7 & 5.93e-02 & 0.5879 & 0.0153 & 0.9847 \\
3 & 0.85 & 6.35e-02 & 0.8126 & 0.0032 & 0.9968 \\
3 & 1 & 8.58e-02 & 1.0335 & 0.0009 & 0.9991 \\
10 & 0.001 & 3.55e-10 & 0.0003 & 0.4994 & 0.5006 \\
10 & 0.005 & 2.39e-09 & 0.0017 & 0.4971 & 0.5029 \\
10 & 0.01 & 1.71e-08 & 0.0033 & 0.4942 & 0.5058 \\
10 & 0.05 & 1.98e-06 & 0.0168 & 0.4708 & 0.5292 \\
10 & 0.1 & 1.57e-05 & 0.0341 & 0.4408 & 0.5592 \\
10 & 0.3 & 1.86e-03 & 0.1227 & 0.2982 & 0.7018 \\
10 & 0.5 & 3.34e-02 & 0.3120 & 0.1014 & 0.8986 \\
10 & 0.7 & 6.03e-02 & 0.5927 & 0.0147 & 0.9853 \\
10 & 0.85 & 6.51e-02 & 0.8177 & 0.0031 & 0.9969 \\
10 & 1 & 8.41e-02 & 1.0475 & 0.0008 & 0.9992 \\
30 & 0.001 & 1.47e-09 & 0.0003 & 0.4994 & 0.5006 \\
30 & 0.005 & 3.71e-09 & 0.0017 & 0.4971 & 0.5029 \\
30 & 0.01 & 6.01e-09 & 0.0033 & 0.4942 & 0.5058 \\
30 & 0.05 & 7.00e-07 & 0.0168 & 0.4708 & 0.5292 \\
30 & 0.1 & 5.44e-06 & 0.0341 & 0.4408 & 0.5592 \\
30 & 0.3 & 2.13e-03 & 0.1230 & 0.2978 & 0.7022 \\
30 & 0.5 & 3.40e-02 & 0.3130 & 0.1007 & 0.8993 \\
30 & 0.7 & 6.06e-02 & 0.5941 & 0.0145 & 0.9855 \\
30 & 0.85 & 6.65e-02 & 0.8234 & 0.0029 & 0.9971 \\
30 & 1 & 8.45e-02 & 1.0500 & 0.0007 & 0.9993 \\
100 & 0.001 & 1.21e-10 & 0.0003 & 0.4994 & 0.5006 \\
100 & 0.005 & 3.73e-09 & 0.0017 & 0.4971 & 0.5029 \\
100 & 0.01 & 2.00e-09 & 0.0033 & 0.4942 & 0.5058 \\
100 & 0.05 & 2.10e-07 & 0.0168 & 0.4708 & 0.5292 \\
100 & 0.1 & 5.51e-06 & 0.0341 & 0.4408 & 0.5592 \\
100 & 0.3 & 2.22e-03 & 0.1230 & 0.2977 & 0.7023 \\
100 & 0.5 & 3.42e-02 & 0.3133 & 0.1004 & 0.8996 \\
100 & 0.7 & 6.06e-02 & 0.5946 & 0.0145 & 0.9855 \\
100 & 0.85 & 6.66e-02 & 0.8239 & 0.0029 & 0.9971 \\
100 & 1 & 8.46e-02 & 1.0510 & 0.0007 & 0.9993 \\
300 & 0.001 & 7.69e-13 & 0.0003 & 0.4994 & 0.5006 \\
300 & 0.005 & 2.31e-09 & 0.0017 & 0.4971 & 0.5029 \\
300 & 0.01 & 1.48e-09 & 0.0033 & 0.4942 & 0.5058 \\
300 & 0.05 & 1.53e-07 & 0.0168 & 0.4708 & 0.5292 \\
300 & 0.1 & 6.66e-06 & 0.0341 & 0.4408 & 0.5592 \\
300 & 0.3 & 2.25e-03 & 0.1231 & 0.2976 & 0.7024 \\
300 & 0.5 & 3.43e-02 & 0.3134 & 0.1004 & 0.8996 \\
300 & 0.7 & 6.07e-02 & 0.5947 & 0.0145 & 0.9855 \\
300 & 0.85 & 6.67e-02 & 0.8241 & 0.0029 & 0.9971 \\
300 & 1 & 8.46e-02 & 1.0512 & 0.0007 & 0.9993 \\
1000 & 0.001 & 1.40e-09 & 0.0003 & 0.4994 & 0.5006 \\
1000 & 0.005 & 2.34e-11 & 0.0017 & 0.4971 & 0.5029 \\
1000 & 0.01 & 3.91e-10 & 0.0033 & 0.4942 & 0.5058 \\
1000 & 0.05 & 2.03e-07 & 0.0168 & 0.4708 & 0.5292 \\
1000 & 0.1 & 7.06e-06 & 0.0341 & 0.4408 & 0.5592 \\
1000 & 0.3 & 2.26e-03 & 0.1231 & 0.2976 & 0.7024 \\
1000 & 0.5 & 3.43e-02 & 0.3135 & 0.1003 & 0.8997 \\
1000 & 0.7 & 6.07e-02 & 0.5948 & 0.0144 & 0.9856 \\
1000 & 0.85 & 6.67e-02 & 0.8242 & 0.0029 & 0.9971 \\
1000 & 1 & 8.46e-02 & 1.0513 & 0.0007 & 0.9993 \\
3000 & 0.001 & 7.44e-12 & 0.0003 & 0.4994 & 0.5006 \\
3000 & 0.005 & 7.39e-11 & 0.0017 & 0.4971 & 0.5029 \\
3000 & 0.01 & 1.12e-10 & 0.0033 & 0.4942 & 0.5058 \\
3000 & 0.05 & 2.17e-07 & 0.0168 & 0.4708 & 0.5292 \\
3000 & 0.1 & 7.18e-06 & 0.0341 & 0.4408 & 0.5592 \\
3000 & 0.3 & 2.26e-03 & 0.1231 & 0.2976 & 0.7024 \\
3000 & 0.5 & 3.43e-02 & 0.3135 & 0.1003 & 0.8997 \\
3000 & 0.7 & 6.07e-02 & 0.5948 & 0.0144 & 0.9856 \\
3000 & 0.85 & 6.67e-02 & 0.8242 & 0.0029 & 0.9971 \\
3000 & 1 & 8.46e-02 & 1.0513 & 0.0007 & 0.9993 \\
\bottomrule
\end{longtable}}

\section{The strict-ridge deficit recheck (completed cells)}
{\small\begin{longtable}{rrrrrr}
\toprule
$\gamma$ & $\tau$ & $F_{\rm strict}$ & deficit (R4) & deficit (strict) & ratio \\
\midrule
\endhead
1 & 0.1 & 4.33e-03 & 7.61e-10 & 1.68e-02 & 0.0 \\
1 & 0.2 & 7.33e-03 & 2.06e-07 & 1.11e-02 & 0.0 \\
1 & 0.3 & 1.45e-02 & 6.66e-06 & 7.54e-03 & 0.0 \\
1 & 0.5 & 3.27e-02 & 8.82e-03 & 2.74e-03 & 3.2 \\
1 & 0.7 & 1.49e-01 & 9.70e-03 & 2.07e-03 & 4.7 \\
1 & 1 & 1.16e-01 & 1.14e-02 & 2.07e-03 & 5.5 \\
10 & 0.1 & 9.02e-03 & 6.16e-09 & 1.45e-02 & 0.0 \\
10 & 0.2 & 3.24e-03 & 1.54e-07 & 2.23e-03 & 0.0 \\
10 & 0.3 & 2.49e-02 & 2.15e-06 & 7.72e-03 & 0.0 \\
10 & 0.5 & 7.58e-02 & 7.18e-03 & 6.30e-03 & 1.1 \\
10 & 0.7 & 5.35e-02 & 2.09e-02 & 7.13e-03 & 2.9 \\
10 & 1 & 3.86e-02 & 1.29e-01 & 6.83e-03 & 18.8 \\
30 & 0.1 & 4.57e-03 & 9.19e-10 & 2.96e-02 & 0.0 \\
30 & 0.2 & 3.15e-03 & 3.67e-08 & 2.23e-03 & 0.0 \\
30 & 0.3 & 3.09e-02 & 9.68e-07 & 1.16e-02 & 0.0 \\
30 & 0.5 & 7.20e-02 & 6.98e-03 & 9.39e-03 & 0.7 \\
30 & 0.7 & 2.82e-02 & 1.75e-02 & 8.34e-03 & 2.1 \\
30 & 1 & 6.39e-02 & 2.16e-01 & 7.43e-03 & 29.0 \\
100 & 0.1 & 6.47e-03 & 1.38e-10 & 3.07e-02 & 0.0 \\
100 & 0.2 & 3.37e-03 & 1.32e-08 & 2.28e-03 & 0.0 \\
100 & 0.3 & 1.12e-02 & 6.41e-07 & 5.99e-03 & 0.0 \\
100 & 0.5 & 4.77e-02 & 7.00e-03 & 9.66e-03 & 0.7 \\
100 & 0.7 & 5.53e-02 & 1.77e-02 & 6.95e-03 & 2.5 \\
100 & 1 & 4.72e-02 & 2.55e-01 & 6.50e-03 & 39.2 \\
300 & 0.1 & 6.87e-03 & 3.93e-11 & 1.03e-02 & 0.0 \\
300 & 0.2 & 1.64e-02 & 8.49e-09 & 7.34e-03 & 0.0 \\
300 & 0.3 & 1.03e-02 & 5.58e-07 & 5.25e-03 & 0.0 \\
300 & 0.5 & 4.88e-02 & 6.98e-03 & 4.67e-03 & 1.5 \\
300 & 0.7 & 3.98e-02 & 2.20e-02 & 5.84e-03 & 3.8 \\
300 & 1 & 6.39e-02 & 2.67e-01 & 3.72e-03 & 71.8 \\
1000 & 0.1 & 5.08e-03 & 1.88e-11 & 2.36e-02 & 0.0 \\
1000 & 0.2 & 3.00e-03 & 7.08e-09 & 2.30e-03 & 0.0 \\
1000 & 0.3 & 1.18e-02 & 5.30e-07 & 6.56e-03 & 0.0 \\
1000 & 0.5 & 8.69e-02 & 6.97e-03 & 6.62e-03 & 1.1 \\
1000 & 0.7 & 8.59e-02 & 1.84e-02 & 6.10e-03 & 3.0 \\
1000 & 1 & 4.15e-02 & 2.71e-01 & 5.48e-03 & 49.5 \\
\bottomrule
\end{longtable}}

\chapter{Repository map}

\section{Solver code (\texttt{projects/REZN/solver\_code/highprec\_dd\_qd/})}
\begin{center}\small\begin{tabularx}{\linewidth}{lX}\toprule
file & role \\
\midrule
\texttt{dd\_ops.py}, \texttt{dd\_k3\_ops.py} & double-double primitives; exact CRRA clearing bisection \\
\texttt{dd\_solver.py} & K=2 DD surface solver (the $\tau$-fold study) \\
\texttt{dd\_k3\_solver.py}, \texttt{dd\_k3\_sweep.py} & K=3 kernel-band R4 production (superseded for economics) \\
\texttt{dd\_k3\_strict\_solve.py} & \strict{} Anderson + deficit metrics \\
\texttt{dd\_k3\_compare.py}, \texttt{dd\_k3\_rconverge.py}, \texttt{dd\_k3\_small\_h.py} & Findings 1, 3 diagnostics \\
\texttt{dd\_k3\_lookup\_alts.py}, \texttt{dd\_k3\_optB\_numba.py}, \texttt{dd\_k3\_optB\_refined.py} & Options A--E; B+ production \\
\texttt{dd\_k3\_optD\_adaptive\_p.py} & adaptive $p$-grid \\
\texttt{dd\_k3\_cheby\_h0.py}, \texttt{dd\_k3\_cheby\_lobatto.py} & spectral program \\
\texttt{dd\_k3\_method\_c\_v2*.py} & structural ansatz: solver, grid, ladder, logit-loss, strict variants \\
\texttt{dd\_k3\_mu\_analytic.py} & closed-form 1D lookup under the ansatz \\
\texttt{dd\_k3\_v3\_fingerprint.py}, \texttt{dd\_k3\_v3\_refit.py} & Finding 7 \\
\texttt{dd\_k3\_lpub\_lpriv\_v2.py}, \texttt{dd\_k3\_delta\_svd.py} & Finding 5 \\
\texttt{dd\_k3\_critpts.py} & 3D critical-point detector (Finding 8) \\
\texttt{dd\_k3\_optB\_pieceCheby.py} & per-segment Chebyshev + the internal cusp detector (Findings 10, 13) \\
\texttt{dd\_k3\_optB\_tail.py}, \texttt{dd\_k3\_optB\_exact.py}, \texttt{dd\_k3\_optB\_multielem.py} & Findings 8, 11, 12 \\
\texttt{dd\_k3\_strict\_smoothed*.py}, \texttt{dd\_k3\_stall\_trace.py} & Finding 9 \\
\texttt{dd\_k3\_strict\_fixed\_fallback.py} & the fallback fix (Finding 14) \\
\texttt{dd\_k3\_zero\_h\_pipeline\_v2/v3/v4.py}, \texttt{dd\_k3\_pipeline\_v5\_final.py} & pipeline v2--v5 \\
\texttt{dd\_k3\_v5\_grid.py} & 36-cell certification (Finding 16) \\
\texttt{dd\_k3\_zero\_h\_solver.py} & Finding 17 \\
\texttt{dd\_k3\_logodds.py}, \texttt{dd\_k3\_logodds\_ladder.py} & the final coordinate configuration + 50-cell ladder \\
\texttt{dd\_k3\_logodds\_newton.py}, \texttt{dd\_k3\_logodds\_ift.py} & IFT Newton (Krylov; HYBRD experiments) \\
\bottomrule\end{tabularx}\end{center}

\section{Pre-K=3 foundations (\texttt{projects/REZN/solved\_fixed\_points/})}
\begin{center}\small\begin{tabularx}{\linewidth}{lX}\toprule
directory & contribution \\
\midrule
\texttt{methodology/} & 9pp method paper + runnable reference operator --- start here \\
\texttt{k3\_coarea\_limit/} & the key anchor: joint $h \to 0$ ladder, PR nailed to $10^{-14}$, deficit $\to 0.28$ \\
\texttt{k3\_verify\_coarea/} & the 19\% naive-scan bias, proven on an analytic surface \\
\texttt{k3\_strict\_h0\_exact/} & flint marching-squares \strict{}: continuity of the bare operator; nails toward $10^{-100}$ \\
\texttt{k3\_hfree\_morse/} & the $C^0$-not-$C^1$ proof ($128\times$ derivative blowup at a Morse price) \\
\texttt{k3\_coarea\_2dsweep/} & the definitive $(\gamma, \tau)$ deficit map (31/32 cells at $10^{-12}$) \\
\texttt{k3\_jensen\_gap/} & the closed-form wedge verified on nailed equilibria \\
\texttt{k3\_cara/} & CARA = $\gamma \to \infty$ benchmark; Hellwig closed-form cross-check \\
\texttt{k4\_*} & K=4: FR globally attracting; PR via extreme-$\gamma$ dropout \\
\texttt{highprec\_dd/} & K=2 DD ladder; the $\tau \approx 0.79$ fold \\
\texttt{flint\_*}, \texttt{k3\_plateau\_*} & superseded naive-scan era (kept as the artifact record) \\
\texttt{FINDINGS.md} & the human-readable index of all of the above \\
\bottomrule\end{tabularx}\end{center}

\chapter{Glossary}
\begin{description}\itemsep2pt
  \item[REE] rational-expectations equilibrium: price function such
    that demands conditioned on (own signal, price) clear at every
    signal profile.
  \item[FR / PR] fully / partially revealing price function.
  \item[deficit] $1 - R^2$ of $\operatorname{logit}P$ on the pooled
    statistic (with stated state-measure); 0 under FR.
  \item[co-area integral] $\int_{\{P=p\}} f/|\nabla P|\, d\sigma$ ---
    the correct level-set evidence measure.
  \item[\strict{}] the bare co-area operator: explicit contours +
    $1/|\nabla P|$, no kernel.
  \item[kernel band] Gaussian-smoothed indicator
    $\int K_h(P - p) f\, du$; consistent quadrature of the co-area
    integral as $h \to 0$.
  \item[central path / joint limit] $h \to 0$ together with
    $\Delta x \to 0$, $\Delta x \ll h$ --- the regime where the kernel
    is both consistent and smooth.
  \item[lookup] $\mu(p, u_k)$: Bayes posterior of $v = 1$ given own
    signal and the price level set.
  \item[POU] partition of unity across the two sweep directions of the
    level-set integral.
  \item[Morse-critical price] $p_c = P(u^*)$ with $\nabla P(u^*) = 0$:
    location of the $C^0$-not-$C^1$ kink in the evidence and of the
    lookup cusps.
  \item[$\Fw$] density-weighted residual
    $\max |\Phi(P) - P| \cdot \prod_i \bar f(u_i)$ --- the economically
    meaningful convergence metric on unbounded grids.
  \item[trivial branch] $P \equiv \tfrac12$: an exact fixed point at
    all parameters; the uninformative pooled-prior equilibrium.
  \item[best iterate] the lowest-residual point of a non-converging
    iteration --- reported as such, never as a fixed point.
  \item[DD / QD] double-double (32 digits) / quad-double (64 digits)
    arithmetic.
\end{description}

\chapter{Chronology of the seventeen findings}
\begin{center}\small\begin{tabularx}{\linewidth}{clX}\toprule
\# & verdict & one-line content \\
\midrule
1 & negative & fixed-grid kernel-band Richardson does not extrapolate to $h = 0$ (kink powers) \\
2 & positive & \strict{} re-solve collapses the R4 deficit $\sim 40\times$; ridge is an artifact \\
3 & negative & cross-variant agreement $\ne$ correctness (robust-but-wrong comparator) \\
4 & positive & Option B+ (cube CDF + PCHIP) wins the builder bake-off; median $10^{-5}$ \\
5 & mixed & $(\Lpriv, \Lpub)$ transform absorbs bulk, preserves cusps; $b(u_k) \approx \tau u_k / 10$ \\
6 & positive & adaptive arclength $p$-grid: $-25\%$ max error, free \\
7 & negative & no finite-parameter polynomial correction reaches $10^{-4}$ (high-rank residual) \\
8 & mixed & 3D critical-point detector works; knots fix the median only \\
9 & negative & $\varepsilon$-smoothed \strict{}: no sweet spot (bias vs convergence) \\
10 & breakthrough & per-segment Chebyshev between cusps: machine eps in-segment \\
11 & negative & tail BCs alone do not move the max error \\
12 & mixed & exact cell mass: halves the median, worsens the max \\
13 & correction & cusps of $\mu$ are \emph{slice}-critical values; internal detector is the right one \\
14 & positive & \texttt{fix\_fallback}: monotone, BC-correct cleaned lookup \\
15 & positive & pipeline v5 assembled: the production zero-$h$ lookup \\
16 & certification & 36/36 cells: median $10^{-16}$, 11/11 monotone, 7 s total \\
17 & negative & smooth lookup inside the FP loop does not break the $G$-truncation wall \\
\bottomrule\end{tabularx}\end{center}

\chapter{Extended figure gallery}

This appendix collects the remaining result figures referenced by the
findings but not embedded in the main text, with interpretive captions.
Together with the in-text figures they constitute the complete graphic
record of the investigation.

\section{Kernel-band era and the full ladder}

\begin{figure}[H]\centering
  \includegraphics[width=0.92\linewidth]{fl_figs/fig2_slope_deficit.png}
  \caption{Slope and deficit along the 1000-cell R4 DD ladder. These are
    the quantities later revised by the \strict{} recheck: the
    high-$\tau$ deficit rise is predominantly operator bias.}
\end{figure}

\begin{figure}[H]\centering
  \includegraphics[width=0.92\linewidth]{fl_figs/fig4_Prange.png}
  \caption{Price range along the R4 ladder: the equilibrium sharpens
    with $\tau$ --- qualitatively robust across operators even where the
    deficit numbers are not.}
\end{figure}

\begin{figure}[H]\centering
  \includegraphics[width=0.92\linewidth]{fl_figs/fig6_contours.png}
  \caption{R4 equilibrium contours across the ladder: the early visual
    evidence of the kink structure that the spectral and cusp analyses
    later formalized.}
\end{figure}

\begin{figure}[H]\centering
  \includegraphics[width=0.85\linewidth]{alts_figs/fig4_distribution.png}
  \caption{Error distributions across the five lookup builders: Option
    B+ concentrates its mass two decades left of all kernel-band
    variants.}
\end{figure}

\begin{figure}[H]\centering
  \includegraphics[width=0.92\linewidth]{optB_figs/fig2_sweep.png}
  \caption{Option B+ accuracy across $(\gamma, \tau)$: uniform in
    $\gamma$, mildly degrading at high $\tau$ as the cube CDF's tails
    thin.}
\end{figure}

\begin{figure}[H]\centering
  \includegraphics[width=0.92\linewidth]{optB_figs/fig3_cross.png}
  \caption{Cross-sectional view: B+ tracks \strict{} to plotting
    accuracy at a representative cell.}
\end{figure}

\begin{figure}[H]\centering
  \includegraphics[width=0.85\linewidth]{optD_figs/fig2_avg.png}
  \caption{Adaptive $p$-grid: average max-error by monitor function ---
    arclength wins consistently.}
\end{figure}

\section{Cusp geometry, continued}

\begin{figure}[H]\centering
  \includegraphics[width=0.85\linewidth]{cusps_figs/fig1_grad_critical.png}
  \caption{Critical points on a second slice/regime, complementing the
    main-text figure: the critical manifold moves with $u_k$ ---
    the geometric root of Finding 13 (slice criticality, not 3D
    criticality, governs the lookup cusps).}
\end{figure}

\begin{figure}[H]\centering
  \includegraphics[width=0.85\linewidth]{cusps_figs/fig1_mu_cusps.png}
  \caption{Lookup cusps for the slice of the previous figure.}
\end{figure}

\begin{figure}[H]\centering
  \includegraphics[width=0.85\linewidth]{cusps_figs/fig4_regime.png}
  \caption{The smooth/cuspy regime boundary in $(\gamma, \tau)$:
    interior critical points (and hence lookup cusps) emerge with
    rising $\tau$; below the boundary global spectral methods would
    suffice.}
\end{figure}

\begin{figure}[H]\centering
  \includegraphics[width=0.8\linewidth]{critpts/dd_k3_strict_fp_g100_t0.2000.png}
  \caption{3D critical points at $(\gamma, \tau) = (100, 0.2)$: 40
    points, 16 unique critical values confined to $p \in [0.34, 0.66]$
    --- the narrower cusp band of the low-$\tau$ regime.}
\end{figure}

\begin{figure}[H]\centering
  \includegraphics[width=0.8\linewidth]{critpts/dd_k3_strict_fp_g1000_t1.0000.png}
  \caption{3D critical points at $(\gamma, \tau) = (1000, 1)$: the
    high-$\gamma$, high-$\tau$ corner; 134 points spanning
    $p \in [0.05, 0.96]$.}
\end{figure}

\section{Structural ansatz, continued}

\begin{figure}[H]\centering
  \includegraphics[width=0.9\linewidth]{mcv2_figs/fig3_Prange.png}
  \caption{Method C v2 ladder: price range vs $\tau$ under the additive
    ansatz --- smooth and monotone by construction.}
\end{figure}

\begin{figure}[H]\centering
  \includegraphics[width=0.9\linewidth]{mcv2_figs/fig4_coefs.png}
  \caption{Evolution of the five ansatz coefficients along the
    $\tau$-ladder: smooth trajectories, suitable for warm-starting and
    for interpolating the ansatz across $\tau$.}
\end{figure}

\begin{figure}[H]\centering
  \includegraphics[width=0.9\linewidth]{mcv2_figs/fig5_diag.png}
  \caption{Ansatz diagnostics along the ladder.}
\end{figure}

\begin{figure}[H]\centering
  \includegraphics[width=0.9\linewidth]{mcv2t_figs/fig2_P_slices.png}
  \caption{Ansatz price slices at $\tau = 0.001$: visually
    indistinguishable from the analytic linear-Gaussian limit.}
\end{figure}

\begin{figure}[H]\centering
  \includegraphics[width=0.85\linewidth]{mcv2g_figs/fig4_cross.png}
  \caption{Grid cross-sections: every $\gamma$ row shows the same
    $\tau$-transition --- the factorization
    $\|F\|(\gamma, \tau) \approx \|F\|(\tau)$ of the ansatz residual.}
\end{figure}

\begin{figure}[H]\centering
  \includegraphics[width=0.9\linewidth]{mcv3_figs/fingerprint.png}
  \caption{The original v3 fingerprint (against the R4-era ansatz fit):
    diagonal and dissenter slices of the residual with candidate
    polynomial overlays --- the first evidence that no simple symmetric
    polynomial captures the non-additive structure.}
\end{figure}

\section{The pipeline iterations (v2--v4) and negative controls}

\begin{figure}[H]\centering
  \includegraphics[width=0.9\linewidth]{v2_figs/g100_t1.0_deg6.png}
  \caption{Pipeline v2 at $(100, 1)$, degree 6: machine epsilon
    in-support, fallback plateau out-of-support (pre-tail-extrapolation
    stage).}
\end{figure}

\begin{figure}[H]\centering
  \includegraphics[width=0.9\linewidth]{v3p_figs/g100_t1.0_deg6.png}
  \caption{Pipeline v3: Chebyshev tail extrapolation with Bayesian BCs
    --- boundary-correct but with degree-2 tail overshoot
    (non-monotone), motivating v4's PCHIP tails.}
\end{figure}

\begin{figure}[H]\centering
  \includegraphics[width=0.9\linewidth]{v4_figs/g100_t1.0_deg6.png}
  \caption{Pipeline v4: shape-preserving tails. Monotone where the
    input is monotone --- which exposed that the raw input was not
    (the fallback artifact), leading to Finding 14.}
\end{figure}

\begin{figure}[H]\centering
  \includegraphics[width=0.9\linewidth]{v5_figs/g1000_t1.0_deg6.png}
  \caption{Pipeline v5 at the $(1000, 1)$ corner: the full
    clean-detect-fit-extrapolate chain on the hardest reference cell.}
\end{figure}

\begin{figure}[H]\centering
  \includegraphics[width=0.9\linewidth]{me_figs/g100_t1.0.png}
  \caption{Multi-element OptB+ (Finding 8): median pinned to zero at
    inserted knots, max unmoved --- the diagnosis that the max error
    lives in the tails, not at the bulk cusps.}
\end{figure}

\begin{figure}[H]\centering
  \includegraphics[width=0.9\linewidth]{exact_figs/compare_g100_t1.0000.png}
  \caption{Exact cell-mass integration vs step summation (Finding 12)
    at $(100, 1)$: the in-support median halves; the tails do not
    improve.}
\end{figure}

\begin{figure}[H]\centering
  \includegraphics[width=0.9\linewidth]{tail_figs/mu_curves_g100_t1.0.png}
  \caption{Tail-BC variants (Finding 11): baseline, BC-augmented,
    logit-space, and Beta-CDF lookups overlaid --- the three BC
    variants are indistinguishable from baseline in the bulk, and the
    bulk is where the max error lives.}
\end{figure}

\begin{figure}[H]\centering
  \includegraphics[width=0.9\linewidth]{smooth_figs/stall_trace.png}
  \caption{The stall trace that opened Finding 9: damped Picard
    \emph{diverges from the saved ``fixed point''} at every damping
    tested --- proof that the saved \strict{} object was a best
    iterate.}
\end{figure}

\begin{figure}[H]\centering
  \includegraphics[width=0.85\linewidth]{mono_figs/fig1_orig_vs_mono.png}
  \caption{Monotonicity projection experiment on the R4 fixed points:
    enforcing axis-monotonicity post hoc reshapes the surface
    materially --- corroborating that the R4 violations were not
    cosmetic.}
\end{figure}

\begin{figure}[H]\centering
  \includegraphics[width=0.75\linewidth]{trf_figs/fig1_trf_conv.png}
  \caption{Trust-region convergence of the ansatz fit: quadratic basin
    behavior, hitting the ansatz expressiveness floor in tens of
    iterations --- solver speed was never the v2/v3 bottleneck.}
\end{figure}

\section{The log-odds configuration, continued}

\begin{figure}[H]\centering
  \includegraphics[width=0.9\linewidth]{lo_figs/g100_t0.1000_G15.png}
  \caption{Log-odds configuration at $G = 15$: the weighted residual
    reaches $10^{-4}$ before the (frozen) $\Lpub$ refresh perturbs it
    --- the best surface obtained in the new coordinates at
    $\tau = 0.1$.}
\end{figure}

\begin{figure}[H]\centering
  \includegraphics[width=0.9\linewidth]{lo_figs/g1_t0.1000_G11.png}
  \caption{$\gamma = 1$ in the log-odds configuration: the
    low-risk-aversion cell --- largest Jensen wedge, widest price
    range, and the noisiest convergence of the sweep, consistent with
    the $1/\gamma$ scaling of the curvature term.}
\end{figure}

\begin{figure}[H]\centering
  \includegraphics[width=0.9\linewidth]{lo_figs/g1000_t0.1000_G11.png}
  \caption{$\gamma = 1000$: approaching the CARA limit --- tighter
    price function, cleaner convergence; the wedge shrinks as
    $1/\gamma$.}
\end{figure}


\end{document}
